%% transmon-benchmark.7 — post-v6 evidence update, 2026-10-11.
%% Draft preparation under #1039, corrected in place on PR #1040 for the
%% qubit-mode finding (#1047/#1055), the Palace unknown count, the C_Sigma
%% mesh-convergence study (#1051, PR #1057) and the refined real-device
%% optimization (#1056, PR #1058).
%% See provenance.md and changelog.md. Historical v6 and its critics are unchanged.

\documentclass{anvil-paper}
\usepackage{xurl}

%% --- URL line-breaking: the two AWS-blog references carry very long
%% hyphenated URLs; allow breaks at hyphens and add stretch so the
%% bibliography stays inside the measure (render-gate overfull check).
\expandafter\def\expandafter\UrlBreaks\expandafter{\UrlBreaks\do\-}
\Urlmuskip=0mu plus 1mu

%% --- Draft-phase helper macros -------------------------------------------
%% Figure placeholder: uses the rendered figure if pub-figures has produced
%% it, otherwise a framed placeholder naming the source script.
\newcommand{\anvilfig}[3][0.85]{%
  \IfFileExists{figures/#2}%
    {\includegraphics[width=#1\linewidth]{figures/#2}}%
    {\fbox{\parbox{#1\linewidth}{\centering\vspace{5em}%
      \texttt{figures/#2}\\[0.5em]%
      {\small placeholder --- render via \texttt{figures/src/#3}}%
      \vspace{5em}}}}}

\title{Differentiable transmon design with a tensor-native finite-element
electromagnetics solver: gradient-based optimization of charging energy,
cross-validated against Palace}
% TODO(operator): final title wording is operator-approved per BRIEF.md
% ⭐ REFRAME (2026-07-16); this is the BRIEF's suggested title class.
\author{Robb Walters\thanks{These authors contributed equally.}
  \and Crutcher Dunnavant\thanks{These authors contributed equally.}}
% TODO(operator): affiliations (author order and equal-contribution
% footnote are operator-resolved per BRIEF.md, 2026-07-14).
\date{\today}

\begin{document}
\maketitle

\begin{abstract}
Superconducting-qubit electromagnetic design is an iterative
guess-and-check loop: the workflow's production field solvers expose no
derivatives of design figures-of-merit, so state-of-practice optimizers
iterate time-consuming electromagnetic simulations and inject parameter
updates from separate analytic physics models
\citep{eriksson2025automated}. We demonstrate the missing capability ---
\emph{solver-derived} design gradients from a 3D finite-element
electromagnetics solve --- and use it for gradient-based optimization of a
transmon's charging energy $E_C$. GEODE-FEM is a Rust solver built on the
Burn tensor framework; a discrete-adjoint layer bridges the one stage its
autodiff tape cannot reach (the sparse factorization), yielding a
validated $2\times 2$ sensitivity matrix --- material ($\partial/\partial
\varepsilon$) and geometry ($\partial/\partial X$) derivatives on both the
electrostatic and the driven $H(\mathrm{curl})$ operators --- each
finite-difference-validated with mutation tests (relative errors
$10^{-9}$--$10^{-5}$). Because the capacitance $C = \phi^{\mathsf{T}}
K\phi$ is variationally stationary, its shape derivative collapses to a
Hellmann--Feynman-like explicit term; composing it with the
capacitance-to-$E_C$ chain, Newton iteration drives a capacitor geometry
to a $0.2156$\,GHz target $E_C$, confirmed by an independent fresh solve
to $1.4\times 10^{-15}$. On the real DeviceLayout transmon, a
finite-difference-validated quadratic-element shape gradient on an
adaptively refined mesh drives the island capacitance $C_\Sigma$ to the
$89.9$\,fF design anchor in three accepted steps, and seven held-out
refinement rounds then bound the converged capacitance of that design
rigorously from above (below $89.02$\,fF) and estimate it by
extrapolation at about $88.8$--$88.9$\,fF; we do not certify a
continuum value. A second design, with the target offset by that
measured bias, converges below $90.08$\,fF, at most $0.20\%$ above the
anchor. Its extrapolated band of $89.84$--$89.96$\,fF contains the
anchor by construction, because its refinement series is the first one
rescaled; at slower assumed convergence rates it lands $0.25$--$0.57\%$
below. Only the upper side is proven, and a two-sided bound is future
work. The convergence study is a result in its own right. Under uniform
refinement the capacitance of the unmodified device falls from
$137.7$\,fF on the committed 133k-tetrahedron mesh to below $83.02$\,fF
(rigorous; estimated at $74.14$--$79.62$\,fF). The gap to the
${\sim}90$\,fF anchor that earlier versions called a model difference
was therefore discretization error, and the converged capacitance of
the committed geometry lies at least $7.6\%$ below the anchor, not
above it. Earlier coarse-mesh optimizations lowered $C_\Sigma$ to reach
the anchor; they moved the wrong way and are kept as a cautionary
precursor. Eigenmode/EPR derivatives remain roadmap. The correctness
credential: from the identical input mesh, GEODE-FEM reproduces Palace's
six eigenmodes above $4.5$\,GHz to $\leq 0.033\%$ (worst case $0.032\%$),
and the transmon qubit mode at $3.45$\,GHz to $5\times10^{-5}\,\%$.
Earlier versions of this paper called that mode spurious and the qubit
mode absent; a Palace run with its search target below the mode
corrects both statements. The contribution is a capability,
not a speed race --- on raw performance the measured record shows no
corner where GEODE-FEM beats Palace at scale, and we report it honestly.
Every result traces to a committed artifact; statements that rest on a
scratch probe, on a reading of Palace's source, or on inference are
labelled as such.
\end{abstract}

\section{Introduction}
\label{sec:intro}

The hard problem in superconducting-qubit electromagnetic design is the
\emph{inverse map}: given target Hamiltonian parameters --- a charging
energy $E_C$, an anharmonicity $\alpha$, coupling capacitances --- produce
a physical layout. Today that map is walked by hand: parameterize a
geometry, run a 3D electromagnetic simulation, extract the parameters,
compare, adjust, repeat. The state-of-practice automated workflow makes
the loop's structure explicit: QDesignOptimizer
\citep{eriksson2025automated} iterates time-consuming electromagnetic
simulations of HFSS-class solvers and guides the parameter updates with
\emph{separate}, user-defined analytic physics models --- precisely
because the field solver at the center of the loop produces observables
but no derivatives of those observables with respect to the design. The
gradient the optimizer actually wants, $\partial(\text{figure of
merit})/\partial(\text{geometry, materials})$ through the field solve
itself, is structurally unavailable from the production solvers ---
commercial (HFSS, COMSOL) and open-source (Palace \citep{palace}) alike:
their assembly and factorization pipelines are hand-written kernels with
no derivative path, and none exposes an adjoint of its solve.

The surrounding literature brackets this gap without closing it. At the
\emph{lumped-circuit} level, differentiation is a solved problem:
SQcircuit \citep{rajabzadeh2023analysis} analyzes arbitrary
superconducting circuits, and its gradient-based follow-up
\citep{rajabzadeh2024general} computes eigenvalue and eigenvector
gradients of the circuit Hamiltonian via PyTorch autodiff, optimizing
qubit designs in circuit-parameter space. But the map from geometry and
materials \emph{to} those lumped parameters still runs through
non-differentiable 3D electromagnetic simulation. At the \emph{field}
level, differentiable electromagnetics is owned by photonics: adjoint and
autodiff methods on FDTD and integral-equation solvers are a mature
inverse-design industry \citep{molesky2018inverse,hughes2019forward,
hammond2022high,schubert2022inverse,mahlau2026fdtdx,ponomareva2025torchgdm},
and differentiable FEM at scale has been proven outside electromagnetics
by JAX-FEM in solid mechanics \citep{xue2023jax}. What none of these
supplies is the regime superconducting-qubit (and, more broadly, RF)
design lives in: \emph{frequency-domain finite elements on body-fitted
tetrahedra with anisotropic materials} --- distributed, field-level EM
gradients from the same FEM solve that produces the design observables.
That wedge is, to our knowledge, uncontested; this paper occupies it.

GEODE-FEM (``geode-fem'') is positioned to close the gap by construction
rather than retrofit. It is a full-wave 3D $H(\mathrm{curl})$/N\'ed\'elec
(plus scalar electrostatic) finite-element solver written in Rust on the
Burn tensor framework \citep{burn}: element assembly is batched
$[n_{\mathrm{elem}},6,6]$ tensor algebra on an autodiff-capable substrate,
so the derivative of assembly with respect to materials and node
coordinates is available where a classical solver has only hand-written
kernels. One stage breaks the tape: the sparse direct factorization
(faer \citep{kazdadi2026faer}) at the heart of every solve. A
discrete-adjoint layer closes exactly that gap --- one extra
(transpose-)solve, reusing the forward factorization, restores the full
gradient (Section~\ref{sec:adjoint}). The resulting capability is
general --- shape and material sensitivities for RF, photonic, and
electrostatic device design --- and the transmon is our demonstration
vehicle because its unmet need is the clearest documented in the
literature, not because the method is qubit-specific.

A design-gradient claim from a new solver is only as credible as the
solver's forward answers, so the paper carries its correctness credential
up front: on a real transmon-plus-readout-resonator geometry meshed once
and given as the same input file to both codes, geode-fem reproduces
Palace --- the MFEM-based \citep{mfem,andrej2024high} reference solver
for superconducting-qubit design --- to within $0.033\%$ across the six
eigenmodes Palace returns above $4.5$\,GHz (worst case $0.032\%$), and to
$5\times10^{-5}\,\%$ on the transmon qubit mode at $3.45$\,GHz, with
analytic scaling tripwires guarding against coincidental agreement
(Section~\ref{sec:agreement}). Palace refines the mesh around the metal
sheets internally, so the two discrete problems are not identical
(Section~\ref{sec:geometry}). This parity is \emph{a credential, not the
contribution} --- matching Palace is the price of entry --- and it is,
incidentally, the first independent third-party validation of the
Palace/DeviceLayout.jl transmon stack, which has no publication of its
own.

Our contributions are:
\begin{itemize}
  \item \textbf{A validated $2\times 2$ discrete-adjoint sensitivity
    matrix for FEM electromagnetics.} Material
    ($\partial/\partial\varepsilon$) and geometry ($\partial/\partial X$,
    exact $\partial K/\partial X$ via forward-mode duals through the
    element kernels) derivatives, on both the scalar electrostatic and
    the complex driven $H(\mathrm{curl})$ N\'ed\'elec operators --- each
    one forward plus one adjoint solve sharing a single factorization,
    each validated against full-pipeline central finite differences with
    mutation tests that reject sign, conjugation, and dropped-term errors
    (Section~\ref{sec:adjoint}).
  \item \textbf{A differentiable capacitance$\to E_C$ chain with a
    variational-stationarity simplification.} Because $C =
    \phi^{\mathsf{T}} K\phi$ is stationary in the solved potential, the
    implicit term vanishes and the shape derivative is a pure
    explicit-geometry (Hellmann--Feynman-like) term; validated against
    the analytic parallel-plate law and central finite differences
    (Section~\ref{sec:chain}).
  \item \textbf{A refinement-validated real-device optimization, with
    its error budget stated.} An end-to-end Newton loop hits a
    $0.2156$\,GHz target $E_C$ on an analytic fixture and is confirmed
    by an independent fresh solve at $1.4\times 10^{-15}$ (the
    parametrization is affine, so one-step convergence proves the loop,
    not a hard optimization). On the real device, a three-parameter
    optimization on an adaptively refined 448k-tetrahedron mesh with
    quadratic elements reaches the $89.9$\,fF anchor in three accepted
    steps, with gradients finite-difference-validated at
    $3.87\times10^{-7}$ or better. Seven held-out refinement rounds
    give a rigorous upper bound on the converged capacitance of the
    design (below $89.02$\,fF) and an extrapolated estimate (about
    $88.8$--$88.9$\,fF). A resolution-corrected second design converges
    below $90.08$\,fF. No continuum value is certified and no lower
    bound is proven (Section~\ref{sec:refined}).
  \item \textbf{A measured mesh-convergence correction.} The committed
    133k-tetrahedron mesh overstates $C_\Sigma$ by at least
    $54.69$\,fF. Converged, the unmodified device lies below
    $83.02$\,fF: under the ${\sim}90$\,fF anchor, where earlier versions
    placed it above. Their ``model difference'' diagnosis is withdrawn,
    and the coarse-mesh optimizations that lowered $C_\Sigma$ toward the
    anchor are re-read as a cautionary precursor
    (Sections~\ref{sec:lom} and~\ref{sec:multiparam}).
  \item \textbf{The correctness credential.} Same-input-mesh,
    same-physics, same-junction-model cross-validation against Palace at
    $\leq 0.033\%$ across the six modes above $4.5$\,GHz and
    $5\times10^{-5}\,\%$ on the qubit mode, with oracle-first methodology
    and honest-physics reporting --- including the correction of two
    claims made in earlier versions, that the qubit mode was absent and
    that the $3.45$\,GHz mode was spurious
    (Sections~\ref{sec:agreement} and~\ref{sec:honest-physics}) ---
    doubling as the first independent validation of the
    Palace/DeviceLayout.jl stack.
  \item \textbf{Mode-verified performance context.} Same-box CPU
    measurements disclose mode selection, field-output and error-estimator
    costs; they replace the earlier unsupported one-core-versus-eight-rank
    efficiency claim. CPU-f64 AMS reduces driven-solve time relative to
    Jacobi, but its iteration and setup costs still grow with refinement
    (Section~\ref{sec:perf}). Eigenmode/EPR derivatives remain roadmap
    (Section~\ref{sec:roadmap}).
\end{itemize}

Section~\ref{sec:related} positions the work; Section~\ref{sec:architecture}
presents the substrate and where its tape breaks;
Sections~\ref{sec:problem}--\ref{sec:honest-physics} establish the
benchmark problem and the correctness credential;
Sections~\ref{sec:lom}--\ref{sec:optdemo} are the contribution ---
the forward LOM chain with its mesh-convergence study, the adjoint
layer, and the optimization tiers; Section~\ref{sec:roadmap} is the eigenmode roadmap;
Sections~\ref{sec:perf}--\ref{sec:discussion} report performance context,
reproducibility, and discussion.

\section{Related Work}
\label{sec:related}

\paragraph{Superconducting-qubit design optimization.}
The design loop this paper attacks is documented across the qubit-design
software stack. QDesignOptimizer \citep{eriksson2025automated} is the
state of practice: physics-guided multi-parameter optimization that
wraps non-differentiable EM simulation (HFSS with pyEPR/Qiskit-Metal-class
post-processing) and steers it with analytic models, because no gradient
is available from the solver. The SQcircuit line
\citep{rajabzadeh2023analysis,rajabzadeh2024general} differentiates at
the lumped-circuit level --- including eigenvalue/eigenvector gradients
of the circuit Hamiltonian via PyTorch --- and optimizes in circuit-parameter
space; the geometry-to-parameters map remains simulation-bound. The
SQuADDS database \citep{shanto2024squadds} attacks the same loop by
lookup: pre-simulated, experimentally-validated design points. Qiskit
Metal's quasi-lumped methods \citep{minev2021circuit} and the
EPR/black-box quantization physics
\citep{koch2007,nigg2012,minev2021} frame the parameter-extraction layer
this paper's LOM chain implements. None of these obtains gradients from
the field solve itself; that is the gap Section~\ref{sec:adjoint} closes.

\paragraph{Differentiable electromagnetics --- owned by photonics.}
Adjoint/autodiff EM is a mature inverse-design discipline in
nanophotonics \citep{molesky2018inverse}, expressed through structured
grids and integral methods: autograd FDFD/FDTD in ceviche
\citep{hughes2019forward}, foundry-constrained JAX inverse design
\citep{schubert2022inverse}, the Meep adjoint pipeline
\citep{hammond2022high}, the JAX-native FDTDX framework
\citep{mahlau2026fdtdx}, GPU-accelerated autodiff integral-method
scattering in TorchGDM \citep{ponomareva2025torchgdm}, and
differentiable RCWA \citep{kim2024meent}. These codes accept rectilinear
grids or Green's-function discretizations; the superconducting-qubit
regime --- body-fitted tetrahedra, anisotropic sapphire,
$H(\mathrm{curl})$ conformity, frequency domain --- is exactly where that
literature stops. Differentiable FEM has been proven outside EM: JAX-FEM
\citep{xue2023jax} (GPU-accelerated, autodiff-native continuum mechanics)
and PyTorch-FEA \citep{liang2023pytorch} in biomechanics; concurrent work
on general differentiable Galerkin assembly (TensorGalerkin
\citep{wen2026learning}, with the torch-sla adjoint sparse-solver library
\citep{chi2026torch}) independently confirms the substrate-level thesis
while reporting nodal magnetostatics rather than full-wave
$H(\mathrm{curl})$ EM. To our knowledge, no prior work delivers
field-level design gradients from a frequency-domain 3D FEM
electromagnetics solve of the class RF/superconducting design uses.

\paragraph{FEM kernels by code generation, and tensor compilers.}
Generating FEM kernels from high-level descriptions is a two-decade
tradition --- FFC \citep{kirby2006compiler}, UFL
\citep{alnaes2014unified}, Firedrake \citep{rathgeber2017firedrake}, TSFC
\citep{homolya2018tsfc} --- whose modern GPU expression is libCEED
\citep{brown2021libceed} and MFEM's partial-assembly path
\citep{andrej2024high}. Palace itself runs on that domain-specific stack.
geode-fem's substrate bet is categorically different: outsource kernels,
scheduling, and backend portability to a \emph{general-purpose ML tensor
stack} (the Halide/TVM compiler lineage
\citep{ragankelley2013halide,chen2018tvm}, reaching scientific computing
via ML frameworks), and phrase assembly and matrix-free application as
batched tensor algebra. For this paper the decisive property of that
choice is not portability but \emph{differentiability}: an
autodiff-capable substrate makes the assembly derivative available by
construction, which is what the adjoint layer of
Section~\ref{sec:adjoint} composes with. Adjacent genera --- simulation
DSLs \citep{hu2019difftaichi,holl2020learning}, neural-ansatz
eigensolvers \citep{jiang2024fieldtnn}, hand-written CUDA FEM
\citep{huantingmeng2014gpu} --- trade away either the general-purpose
stack or the discretization itself.

\paragraph{Validation and cross-solver benchmarking.}
Code-to-code comparison on precisely specified problems is an established
genre in computational electromagnetics (the TEAM workshop tradition
since 1986\footnote{The TEAM founding documents are workshop reports
without resolvable DOIs; we describe the tradition in prose.}; nano-optics
solver comparisons \citep{hoffmann2009comparison}); in V\&V vocabulary
\citep{roache1997quantification,oberkampf2010verification} our credential
is verification-side, code-to-code evidence. On the qubit-EM axis the
closest works are SQDMetal \citep{sommers2025open} (Palace cross-checked
against COMSOL/HFSS at the $0.02$--$0.13\%$ level, independently
constructed models) and the Palace-centered workflow of
\citet{ye2025electromagnetic} (validated against cryogenic measurement at
${\sim}0.3\%$). Both wrap and trust Palace; neither is a same-mesh
comparison against an independently \emph{implemented} solver. Our
protocol removes mesh generation and model entry as confounders ---
identical MSH fixture, identical lumped-shunt junction, first-order
N\'ed\'elec on both sides --- so residuals reflect formulation, solver
and Palace's internal refinement of the metal sheets
(Section~\ref{sec:geometry}), not independently built models.
Mature open EM tools normally carry a methods paper (GetDP
\citep{dular1998general}, FEniCS/DOLFIN \citep{logg2010dolfin}, Gmsh
\citep{gmsh}, scikit-rf \citep{arsenovic2022scikit}, MFEM
\citep{mfem}); Palace and DeviceLayout.jl are unpublished outliers
\citep{palace,palace-blog,devicelayout,devicelayout-blog,
devicelayout-juliacon}, which is the gap the credential narrows from the
outside. geode-fem's numerics: first-order N\'ed\'elec elements
\citep{nedelec1980}, shift-invert Lanczos \citep{arpack} over a real
symmetric pencil factorized with faer \citep{kazdadi2026faer}; Palace:
SLEPc-class Krylov \citep{slepc} with divergence-free projection.
The timed Palace build resolves its default linear solver to SuperLU
under GMRES; explicit AMS runs are reported separately
\citep{hiptmair-xu,tzaniovkolevpanayotsvassilevski2018parallel}; the gauge-fixing literature
behind Section~\ref{sec:honest-physics} is the tree--cotree tradition
\citep{albanese1988solution,manges1995generalized}.

\section{GEODE-FEM: A Tensor-Native Substrate, and Where the Tape Breaks}
\label{sec:architecture}

geode-fem is written in Rust on the Burn tensor framework \citep{burn},
whose cubecl JIT compiles tensor kernels for CPU, CUDA, and WebGPU
backends from one source. Every throughput-critical FEM stage is phrased
as the batched dense algebra the stack already optimizes.

\paragraph{Batched element-local assembly.}
For first-order N\'ed\'elec tetrahedra the element curl--curl and mass
contributions are $6\times 6$ matrices depending on per-element geometry
factors and material tensors (including the rotated anisotropic sapphire
permittivity of Section~\ref{sec:geometry}); the scalar electrostatic
operator is the analogous P1 kernel. geode-fem computes them for all
elements at once as batched $[n_{\mathrm{elem}},6,6]$ tensor operations
--- one kernel launch per assembly stage rather than an element loop ---
then either scatters into a global sparse matrix (the assembled f64 CPU
path) or retains the batch for matrix-free application.

\paragraph{Matrix-free application and on-device Krylov.}
The matrix-free operator apply is a three-stage tensor pipeline ---
gather global DOFs to element-local vectors, batched matmul against the
stored element operators, scatter-add back --- and the driven
frequency-sweep solve runs a conjugate orthogonal conjugate gradient
(COCG) iteration entirely on the device with a constant number of scalar
host--device synchronizations per iteration. The backend (ndarray on CPU,
CUDA on GPU) is a type parameter, not a code path.
Section~\ref{sec:perf} reports honestly how far below the
GPU-amortization threshold the current benchmark sizes sit.

\paragraph{Honest constraints.}
Two boundaries are constraints, not footnotes. The CUDA backend
(burn-cuda 0.21) is f32-only today --- the underlying cubecl runtime
currently disables f64\footnote{Upstream constraint of the burn/cubecl
stack at version 0.21 \citep{burn}; tracking-issue URL to be added at
revision. TODO(operator).} --- so every GPU result is
mixed-precision-qualified. And the eigensolve that produces the
credential comparison is factorization-bound: its inner solves are faer
sparse-LU factorizations \citep{kazdadi2026faer}, a classical serial CPU
component outside the tensor-compiler story.

\paragraph{Where the tape breaks --- and why an adjoint layer, not
autodiff-through-the-solver.}
The substrate choice has a consequence the rest of this paper builds on:
Burn's automatic differentiation reaches everything \emph{up to} the
linear solve --- material and node-coordinate dependence of the assembled
operators is differentiable by construction --- but the sparse
factorization is an external, opaque kernel that breaks the tape.
Differentiating \emph{through} a factorization by unrolling is neither
practical nor necessary: the classical discrete-adjoint identity turns
the missing derivative into one additional linear solve against the
\emph{same} factorization (Section~\ref{sec:adjoint}). The division of
labor is deliberate --- autodiff where the substrate provides it (element
kernels, via forward-mode duals), adjoint algebra where it does not (the
solve) --- and it is what makes the gradients in this paper
\emph{solver-derived} in the sense the design loop of
Section~\ref{sec:intro} lacks: the same assembly and the same
factorization that produce the observable produce its gradient.

\section{Benchmark Problem: Geometry, Mesh, and Junction Model}
\label{sec:problem}

\subsection{Geometry and mesh provenance}
\label{sec:geometry}

The device is the \texttt{SingleTransmon} example shipped with
DeviceLayout.jl v1.15.0 \citep{devicelayout}: a transmon with readout
resonator on a sapphire substrate, enclosed in a vacuum box. The sapphire
is anisotropic, with permittivity tensor
$\varepsilon = R\,\mathrm{diag}(9.3,\, 9.3,\, 11.5)\,R^{\mathsf{T}}$
rotated approximately $36.87^\circ$ in-plane. The DeviceLayout.jl
workflow emits seven named physical groups (metal, substrate, vacuum,
junction lumped-element surface, two readout ports, exterior boundary)
and drives Gmsh \citep{gmsh} to produce a tetrahedral mesh, exported as
an MSH~4.1 fixture and pinned by SHA-256 (Section~\ref{sec:repro}).

The mesh has 22{,}684 nodes and 133{,}314 tetrahedra, yielding 156{,}863
N\'ed\'elec edge unknowns, of which 133{,}108 interior degrees of freedom
remain after eliminating perfect-electric-conductor (PEC) boundary edges.
PEC is applied on all metal surfaces and the exterior box; the two
readout ports are left open (lossless), so no resistive element enters
and the eigenproblem pencil stays real symmetric.
These counts are geode-fem's. Palace reads the same file but does not
solve the same discrete problem: its run logs record $4{,}913$ elements
added by local bisection around under-resolved interior boundaries and
$6{,}014$ duplicated vertices where it cracks the metal sheets, and it
assembles $183{,}891$ first-order N\'ed\'elec unknowns against
geode-fem's $156{,}863$ (about $17\%$ more). ``Same mesh'' in this paper
therefore means the same \emph{input} mesh. Earlier versions said the
two solvers consumed an identical mesh and solved an identical discrete
problem; that was wrong, and the timing comparisons of
Section~\ref{sec:perf-cpu} carry the difference as a fairness caveat.
This mesh is also coarse for the electrostatic quantities: one uniform
refinement lowers $C_\Sigma$ by $26.3\%$ (Section~\ref{sec:lom}).
Figure~\ref{fig:geometry} shows the geometry with physical groups
color-coded. Every measurement in this paper starts from this one
committed mesh. The eigenmode credential and the coarse-mesh
electrostatic chain run on it directly; the scale study of
Section~\ref{sec:perf} uses a uniformly refined variant; and the
convergence study and refined optimization of Sections~\ref{sec:lom}
and~\ref{sec:refined} use uniform and adaptive refinements of it that
leave its polygonal surfaces unchanged.

\begin{figure}[t]
  \centering
  \anvilfig{fig1-geometry.pdf}{fig1\_geometry.py}
  \caption{Benchmark geometry: top-down projection of the committed
  surface mesh (the SHA-256-pinned MSH~4.1 fixture) for the
  DeviceLayout.jl v1.15.0 \texttt{SingleTransmon} transmon and readout
  resonator. Surface physical groups are color-coded --- the metal ground
  plane (with the coplanar-waveguide feedline, meander resonator, and
  transmon slot as etched gaps), the four-triangle junction
  \texttt{lumped\_element} surface (enlarged in the inset, actual
  triangulation shown), the two readout ports, and the exterior-boundary
  outline; the \texttt{substrate} and \texttt{vacuum} volume groups fill
  the enclosing box. Both solvers read this file (22{,}684
  nodes; 133{,}314 tetrahedra); Palace then refines it internally
  around the metal sheets (Section~\ref{sec:geometry}).}
  \label{fig:geometry}
\end{figure}

\subsection{The junction model: a lumped reactive shunt}
\label{sec:junction}

The Josephson junction is modeled identically in both solvers as a lumped
reactive shunt --- a parallel $L \parallel C$ element with
$L = 14.860$\,nH and $C = 5.5$\,fF --- attached to the four-triangle
\texttt{lumped\_element} surface group $\Gamma$. Following the uniform
lumped-port construction, with port length $\ell$ and width $w$, the shunt
contributes
\begin{equation}
  K_{\mathrm{port}} = \frac{\ell}{w\,L}\, S_\Gamma
  \qquad\text{and}\qquad
  M_{\mathrm{port}} = \frac{C\,\ell}{w}\, S_\Gamma ,
  \label{eq:port}
\end{equation}
where $S_\Gamma$ is the port-surface matrix assembled on $\Gamma$. The
eigenproblem is then the real symmetric generalized pencil
\begin{equation}
  \left(K + K_{\mathrm{port}}\right) x
  \;=\; \omega^2 \left(M + M_{\mathrm{port}}\right) x .
  \label{eq:pencil}
\end{equation}
In Palace, the same physics is entered as a reactive \texttt{LumpedPort}
($L \parallel C$, no resistance) on the same surface group; the readout
ports are configured open, so Palace's pencil is likewise real. Palace's
source confirms the correspondence: its \texttt{LumpedPort} adds the
same scalar surface-mass terms with the same per-square scalings
(\path{docs/research/transmon-lumped-port-formulation.md}, \S1). This
one-to-one correspondence is what makes the comparison same-physics as
well as same-input-mesh. It also means that any property of the
distributed surface term in Eq.~\ref{eq:port} is shared by both solvers
(Section~\ref{sec:port-artifact}).

\section{The Correctness Credential: Same-Input-Mesh Cross-Validation
Against Palace}
\label{sec:agreement}

A solver claiming design gradients must first prove its forward answers;
this section is that proof. Palace's own agreement with commercial tools
is established \citep{sommers2025open}, so reproducing Palace is the bar
a new implementation must clear before anything it computes \emph{about}
its solutions (such as their derivatives) is worth discussing.

\paragraph{Protocol.}
The benchmark is oracle-first. The mesh is generated once, committed,
pinned by SHA-256, and read by both solvers at matched (first) order
with the identical junction model; residual disagreement is therefore
attributable to formulation, implementation or Palace's internal
refinement (Section~\ref{sec:geometry}), not to independent meshing.
Two analytic tripwires guard against coincidental or
self-consistent-but-wrong agreement. Both test the port term's $L$ and
$C$ scalings. The first is the resonance of the port sheet's own
inductance against its own capacitance,
$f_{LC} = 1/(2\pi\sqrt{LC}) = 17.60$\,GHz, which anchors the
$17.49$\,GHz mode. The second is a $\sqrt{L}$-scaling test: doubling $L$
must move that mode by $1/\sqrt{2}$ --- measured: $17.4901 \to
12.37$\,GHz, committed ratio $0.7071$ (computed from the unrounded
values; Figure~\ref{fig:lscaling}). Earlier versions called the
$17.49$\,GHz mode ``the junction mode'' and treated these two checks as
pinning the junction physics. Section~\ref{sec:port-artifact} shows that
this mode carries no port voltage, so the checks pin the port
\emph{operator}, not the qubit; the qubit mode has its own scaling check
in Section~\ref{sec:qubit-mode}. The suite also contains inverse tests
that \emph{must fail} (perturbed configurations required to disagree with
the oracle --- a pipeline that still passes under them is not measuring
anything), and the Palace oracle is deterministic: a rerun from the
committed configuration reproduced its committed \texttt{eig.csv}
bit-for-bit.

\paragraph{Mode identification.}
Modes are identified, not eyeballed. Within geode-fem, the
stiffness-participation ratio
\begin{equation}
  p \;=\; \frac{x^{\mathsf{T}} K_{\mathrm{port}}\, x}
               {x^{\mathsf{T}} \left(K + K_{\mathrm{port}}\right) x}
  \label{eq:participation}
\end{equation}
measures the share of a mode's stiffness energy stored in the port
surface term, whatever its spatial pattern. It is $p = 1.000$ for the
$17.49$\,GHz mode, $p = 0.9942$ for the $3.45$\,GHz mode, and $p = 0.000$
(at the artifact's three-decimal precision) for the resonator and box
modes. Palace's port EPR is built from the port \emph{voltage}
functional alone, so it counts only the uniform-voltage part of that
energy. The two diagnostics agree where the port field is a uniform
voltage and disagree where it is not: on the qubit mode Palace reports
$0.99424$ against geode-fem's $0.9942$, while on the $17.49$\,GHz mode
Palace reports $2.49\times10^{-8}$ against geode-fem's $1.000$
(Section~\ref{sec:honest-physics}). In the original Palace run, which
searched only above $4.5$\,GHz, all six port-EPR values are
$\leq 5\times10^{-4}$; earlier versions read that as two diagnostics
ranking the modes differently, when the cause was the qubit mode's
absence from Palace's search window. The six-mode comparison of
Table~\ref{tab:agreement} is matched by frequency.

\begin{table}[t]
  \centering
  \caption{Eigenmode agreement from the identical input mesh (22{,}684
  nodes / 133{,}314 tets; geode-fem solves 156{,}863 N\'ed\'elec DOFs,
  133{,}108 interior after PEC; Palace refines internally and solves
  183{,}891), first order in both solvers, junction as a lumped reactive
  shunt ($L = 14.860$\,nH, $C = 5.5$\,fF). These are the six modes
  Palace returns when it searches above $4.5$\,GHz; the qubit mode below
  that target is in Table~\ref{tab:qubit}. geode-fem at commit
  \texttt{3174015}; Palace at commit \texttt{fba6a5b}. All values are
  quoted verbatim from the committed
  \texttt{benchmarks/transmon\_eigen/results.toml}: geode-fem frequencies
  at the committed three decimals, Palace at full precision, and the
  committed per-mode \texttt{rel\_err\_pct}, computed against the rounded
  geode-fem values (precision disclosure in Section~\ref{sec:repro}).}
  \label{tab:agreement}
  \small
  \begin{tabular}{@{}llccc@{}}
    \toprule
    Mode & & geode-fem (GHz) & Palace (GHz) & $\Delta$ (\%) \\
    \midrule
    \texttt{resonator} & readout resonator & 5.153 & 5.151335830348 & 0.032 \\
    \texttt{mode\_2} & box/cavity & 15.465 & 15.46052107794 & 0.029 \\
    \texttt{junction\_lc} & port-sheet LC$^{\ast}$ & 17.490 & 17.49010903536 & 0.001 \\
    \texttt{mode\_4} & box/cavity & 18.693 & 18.69165792915 & 0.007 \\
    \texttt{mode\_5} & box/cavity & 20.703 & 20.69755679425 & 0.026 \\
    \texttt{mode\_6} & box/cavity & 26.088 & 26.08089940472 & 0.027 \\
    \bottomrule
  \end{tabular}

  \smallskip
  {\footnotesize $^{\ast}$The artifact key \texttt{junction\_lc} is kept
  verbatim. Palace's port EPR for this mode is $2.49\times10^{-8}$, and
  Section~\ref{sec:port-artifact} identifies it as an artifact of the
  distributed port term.}
\end{table}

\begin{figure}[t]
  \centering
  \anvilfig{fig2-agreement.pdf}{fig2\_agreement.py}
  \caption{Mode-frequency agreement: geode-fem versus Palace from the
  identical input mesh, for the six modes above Palace's $4.5$\,GHz
  search target, with per-mode relative deviations annotated. Data from
  the committed benchmark artifacts
  (\texttt{benchmarks/transmon\_eigen/results.toml} and Palace's
  committed \texttt{eig.csv}).}
  \label{fig:agreement}
\end{figure}

\begin{figure}[t]
  \centering
  \anvilfig{fig3-lscaling.pdf}{fig3\_l\_scaling.py}
  \caption{Inductance-scaling checks on log-log axes. Doubling $L$ moves
  the $17.49$\,GHz port-sheet mode from $17.4901$ to $12.37$\,GHz
  (committed ratio $0.7071 = 1/\sqrt{2}$), and moves the qubit mode from
  $3.4528$ to $2.4449$\,GHz in geode-fem (ratio $0.7081$) and from
  $3.452749$ to $2.444899$\,GHz in Palace (ratio $0.70810$). The qubit
  mode's ratio sits slightly above $1/\sqrt{2}$ because the full-wave
  mode lies $0.28\%$ below the lumped estimate at $L$ and $0.14\%$ below
  it at $2L$. Values from the committed
  \texttt{benchmarks/transmon\_eigen/results.toml}
  (\texttt{[tripwires.junction\_l\_doubling]}, \texttt{[spurious\_mode]}
  and \texttt{[oracles.palace\_target2]}).}
  \label{fig:lscaling}
\end{figure}

\paragraph{Result.}
Table~\ref{tab:agreement}: all six modes agree to $\leq 0.033\%$ (worst
case $0.032\%$). The $17.49$\,GHz port-sheet mode agrees most closely,
at $0.001\%$, and sits $0.6\%$ below $f_{LC} = 17.60$\,GHz in \emph{both}
solvers. Agreement that tight on a mode confined to the four port
triangles is what an identical port operator would produce, and it is
weaker evidence about the bulk discretizations than the cavity modes'
${\sim}0.03\%$. The qubit mode, which Palace returns only when asked to
search below it, agrees to $5\times10^{-5}\,\%$
(Section~\ref{sec:qubit-mode}, Table~\ref{tab:qubit}). The agreement is
between genuinely independent code paths --- a batched-tensor assembly
feeding a direct-factorization Lanczos in Rust, and a projected,
preconditioned Krylov method in C++ --- on discrete problems that
differ by Palace's internal refinement ($183{,}891$ against $156{,}863$
unknowns). A common implementation error is therefore an unlikely
explanation for the residual. A shared \emph{modeling} property is not
excluded by cross-comparison, and Section~\ref{sec:port-artifact}
reports one. All of this is same-input-mesh agreement at first order;
none of these frequencies is a mesh-converged value.

\section{Honest Physics: The Qubit Mode, Two Withdrawn Claims, and a
Port-Sheet Artifact}
\label{sec:honest-physics}

Versions 1--6 of this paper made two claims in this section, and both
were wrong. The first was that no qubit mode exists in this junction
model by construction, and that Palace confirmed its absence. The second
was that geode-fem's $3.4528$\,GHz mode is a spurious port artifact with
no Palace counterpart. We withdraw both. The committed evidence
identifies the $3.45$\,GHz mode as the transmon qubit mode, which Palace
also returns, and it points to a different mode, at $17.49$\,GHz, as the
artifact of the port model. This section separates what is proven by
committed artifacts from what is derived, what a scratch probe measured,
and what is only planned.

\subsection{The \texorpdfstring{$3.45$\,GHz}{3.45 GHz} mode is the transmon qubit mode}
\label{sec:qubit-mode}

\paragraph{Why the earlier Palace run could not return it.}
The original Palace run asked for six modes with
\texttt{Target}~$= 4.5$\,GHz and returned nothing below the $5.15$\,GHz
resonator. Earlier versions read that as Palace hunting near
$4.5$\,GHz and finding no qubit mode. That misreads \texttt{Target}.
Palace's eigensolver accepts only eigenvalues at or above the target: its
SLEPc configuration restricts the search region to that side, and its
configuration schema defines \texttt{Target} as the frequency ``above
which to search''
(\path{docs/research/transmon-lumped-port-formulation.md}, \S2, read
from the Palace source at the benchmarked changeset). The run's silence
below $4.5$\,GHz was therefore not evidence of absence.

\paragraph{Proven: Palace returns the mode when asked.}
We re-ran Palace at the same changeset on the same input mesh, with the
same materials, \texttt{LumpedPort}, order, eigensolver tolerances and
linear solver, changing only \texttt{Target} ($4.5 \to 2.0$\,GHz) and the
number of modes ($6 \to 8$); a second cell also doubles $L$. A control
cell with the original request reproduces the committed spectrum to a
relative $6.5\times10^{-13}$, so the local build does not move the
eigenvalues. Table~\ref{tab:qubit} gives the result
(\path{benchmarks/transmon_eigen/results.toml},
\texttt{[oracles.palace\_target2]}; raw outputs under
\path{reference/fixtures/transmon_palace/results_p1_target2/}).

\begin{table}[t]
  \centering
  \caption{The qubit mode in both solvers, from the identical input mesh
  at first order. Palace values are from the \texttt{Target}~$= 2$\,GHz
  re-run; geode-fem values are the committed ones that earlier versions
  reported as a spurious mode. The lumped estimate is
  $1/(2\pi\sqrt{L(C_\Sigma + C_J)})$ with the electrostatic
  $C_\Sigma = 136.684731$\,fF of Section~\ref{sec:lom} and
  $C_J = 5.5$\,fF. Palace's column is its port EPR; geode-fem's is the
  stiffness participation of Eq.~\ref{eq:participation}. This is a
  same-input-mesh comparison on a coarse mesh, not a converged qubit
  frequency.}
  \label{tab:qubit}
  \small
  \begin{tabular}{@{}lccc@{}}
    \toprule
    Quantity & Palace & geode-fem & Lumped estimate \\
    \midrule
    Qubit mode, $L = 14.860$\,nH (GHz) & 3.452749 & 3.452751 & 3.46246 \\
    Port EPR / participation $p$ & 0.99424 & 0.9942 & --- \\
    Qubit mode, $2L$ (GHz) & 2.444899 & 2.4449 & 2.44833 \\
    $L$-doubling ratio ($1/\sqrt{2} = 0.70711$) & 0.70810 & 0.7081 & 0.70711 \\
    \midrule
    $17.49$\,GHz mode: port EPR / $p$ & $2.49\times10^{-8}$ & 1.000 & --- \\
    $17.49$\,GHz mode at $2L$ (GHz) & 12.3674 & 12.37 & --- \\
    \bottomrule
  \end{tabular}
\end{table}

Palace finds the mode at $3.452749$\,GHz with port EPR $0.99424$;
geode-fem's committed value is $3.452751$\,GHz with $p = 0.9942$. The two
frequencies agree to $5\times10^{-5}\,\%$. Under $L$ doubling the mode
moves to $2.444899$\,GHz in Palace, a ratio of $0.70810$, against
geode-fem's committed $0.7081$. Palace's value lies $0.280\%$ below the
lumped estimate at $L$ and $0.140\%$ below it at $2L$. That estimate
takes $C_\Sigma$ from the electrostatic branch on the same coarse mesh,
so it is a consistency check between this paper's eigenmode and LOM
branches, not an independent anchor. Every other mode in the re-run
matches the committed Palace spectrum, and two further modes converge
above it ($33.384$ and $33.896$\,GHz).

Earlier versions explained the supposed absence by saying that the
lumped-shunt model gives the junction only its own $5.5$\,fF, so that its
resonance lands at $17.5$\,GHz, and that the pad capacitance of a
physical transmon \citep{koch2007} was missing. That explanation was
wrong. The full-wave eigenproblem contains the island's capacitance
through the field itself, which is why a mode exists in which the
junction inductance resonates against the whole island.

\paragraph{What this does not show.}
This is not the converged physical qubit frequency. Both solvers start
from the same coarse 133k-tetrahedron mesh at first order, and they do
not solve the same discrete problem: Palace refines around the metal
sheets and solves $183{,}891$ unknowns against geode-fem's $156{,}863$
(Section~\ref{sec:geometry}). The mesh is also far from converged for
the capacitance that sets this frequency: one uniform refinement lowers
the scalar-$\varepsilon$ $C_\Sigma$ by $26.3\%$
(Section~\ref{sec:lom}). The lumped estimate goes as
$1/\sqrt{C_\Sigma + C_J}$, so we infer that the qubit frequency is not
converged on this mesh either; that inference has not been tested by a
refined eigensolve. The DeviceLayout.jl blog workflow
\citep{devicelayout-blog} reports a spectrum from $4.14$ to
$5.591$\,GHz that includes a ${\sim}4$\,GHz qubit mode. We do not claim
to reproduce it, and we do not attribute the difference from our
$3.45$\,GHz to any single cause.

\paragraph{The committed artifacts still carry the old labels.}
\path{benchmarks/transmon_eigen/results.toml} keeps this mode under the
section name \texttt{[spurious\_mode]} with \texttt{excluded = true}, and
the CPU timing artifacts label it \texttt{spurious\_port}. The
re-baselining of those classifications is tracked separately (issue
\#1045) and has not been done. We quote the artifact keys verbatim and
state here what the evidence now shows.

\subsection{The gauge and projection arc, re-read}
\label{sec:gauge-arc}

Earlier versions told the following arc as the diagnosis of a spurious
mode. Its measurements stand, every step committed
(Section~\ref{sec:repro}); what they mean changes once the two modes are
correctly identified. geode-fem's headline Lanczos operates on the raw
pencil of Eq.~\ref{eq:pencil} without a gauge projection. (1) Classical
tree--cotree elimination
\citep{albanese1988solution,manges1995generalized} is not
spectrum-preserving for the generalized eigenproblem --- the resonator
shifts $1.64\%$ ($5.1528 \to 5.2372$\,GHz), outside our bar, pinned as a
regression test. (2) The spectrum-preserving bulk divergence-free
projector (structurally Palace's confinement) preserves the cavity
spectrum ($0.029\%$) and collapses the 13{,}747-mode gradient cluster.
It removes the $17.49$\,GHz mode, which is almost entirely a gradient
field (measured divergence ratio $50.2$; projected norm
$1.06\times10^{-4}$), and it keeps the $3.45$\,GHz mode, which is
solenoidal with respect to the interior gradients (divergence ratio
${\approx}6\times10^{-15}$). Earlier versions described this as the
projector deflating the physical junction mode while the spurious mode
survived. With the identities corrected it reads the other way: the
projector kept the qubit mode and removed the port-sheet mode. (3) A
port-aware rank-1 re-admission of one gradient direction restores the
$17.49$\,GHz mode and returns all six of Palace's
\texttt{Target}~$= 4.5$\,GHz modes at $\leq 0.029\%$ worst case. That
construction did what it was built to do. Its purpose, though, was to
retain a mode that Section~\ref{sec:port-artifact} identifies as an
artifact, because the Palace reference then in hand contained that mode
and did not contain the qubit mode.

The remaining committed characterization of the $3.45$\,GHz mode is
unchanged and now reads as a description of the qubit mode: under the
$L$-doubling harness it moves $3.4528 \to 2.4449$\,GHz (ratio $0.7081$)
with $99.4\%$ of its stiffness energy in the $K_{\mathrm{port}}$ term.
Earlier versions added that Palace's \texttt{LumpedPort} represents the
port differently and therefore never sees this mode. That was also
wrong: Palace's port operator has the same structure
(Section~\ref{sec:junction}), and Palace sees the mode.

The headline comparison deliberately remains the committed ungauged
benchmark. The subsequent port-subspace implementation removes the
near-zero survivor and admits the port-patch gradient direction with one shifted
factorization. On this fixture, requesting seven modes at $4.5$\,GHz
returns the six modes of Table~\ref{tab:agreement} (worst relative
frequency error $0.0289\%$) plus the $3.452751$\,GHz qubit mode;
requesting six still misses the highest of the six. The committed local
comparison records mean process CPU times of $46.32$\,s for the older
eight-mode extraction route and $23.25$\,s for the new seven-mode route.
These are geode-only measurements on a loaded Apple M3 Ultra, not a
Palace comparison or a same-box update to the EC2 timings
(Section~\ref{sec:perf-cpu}). geode-fem's benchmark harness still
classifies modes by frequency matching against the Palace reference; no
reference-free rule is built. Figure~\ref{fig:participation} shows the
spectrum.

\subsection{The \texorpdfstring{$17.49$\,GHz}{17.49 GHz} mode: a zero-voltage resonance of the distributed port term}
\label{sec:port-artifact}

\paragraph{Proven by committed data.}
Palace's port EPR for the $17.49$\,GHz mode is $2.49\times10^{-8}$, in
the original run and in the re-run. After $L$ doubling the mode moves to
$12.3674$\,GHz (ratio $0.70711$) and its port EPR is
$2.94\times10^{-8}$. Palace builds that EPR from the port voltage, so in
Palace this mode carries essentially no port voltage. In geode-fem the
same mode has $p = 1.000$: all of its stiffness energy is in the port
surface term. Both statements hold only if that energy sits in a
surface pattern whose net voltage is zero.

\paragraph{Measured by a scratch probe.}
A probe on the committed fixture measured the uniform-voltage share of
each mode's port energy in geode-fem: $0.00\%$ for the $17.49$\,GHz mode
and $100.00\%$ for the $3.45$\,GHz mode. The probe source is committed
(\path{docs/research/transmon-lumped-port-formulation.probe-1003.rs})
but it is scratch evidence: the build does not compile it, and no
regression test pins these two numbers.

\paragraph{Derived.}
The junction patch is four triangles with five nodes, of which only the
centre node is not on a conductor. The gradient of that node's hat
function is curl-free and has zero port voltage, but the isotropic
surface term of Eq.~\ref{eq:port} still stores energy in it. A field
confined to the patch then has Rayleigh quotient $1/(LC)$, which is
$f_{LC} = 17.60$\,GHz, and the volume mass of the adjacent tetrahedra
lowers it to $17.49$\,GHz. A two-terminal junction cannot carry this
pattern. On this reading the mode is a sheet-internal resonance created
by representing a lumped element as a distributed isotropic sheet, and
both solvers have it because both use that representation. The reading
predicts one such mode per free patch node, so a finer port mesh should
show a cluster just below $f_{LC}$. That falsifier has not been run.

\paragraph{What it means for the credential.}
The $0.001\%$ agreement on this mode is agreement on a shared artifact
of the port model. It still tests that the two implementations assemble
the same port operator, and nothing more. The physical content of the
credential is the five cavity modes and the qubit mode.

\paragraph{Roadmap, not built.}
The planned remedy is a rank-1 port-voltage formulation (issue \#1044):
replace $S_\Gamma$ in Eq.~\ref{eq:port} by $b\,b^{\mathsf{T}}$, with
$b_i = (1/w)\int_\Gamma N_i\cdot\hat{e}\,dS$ the port-voltage
functional, so that only the energy of the port voltage is kept. In the
same scratch probe this formulation removes the $17.49$\,GHz mode
(no mode remains between $15.47$ and $18.69$\,GHz), leaves the
$3.4528$\,GHz mode and the five cavity modes unchanged to four decimals,
and makes the plain bulk projector exact (six modes from a six-mode
request, residuals $\leq 1.7\times10^{-9}$, no reference filter). These
are probe measurements on uncommitted solver code. No shipped solver
path, committed benchmark or table in this paper uses the formulation.
A solver using it would no longer reproduce Palace's $17.49$\,GHz mode,
and whether it becomes the default is an open decision (issue \#1046).

\begin{figure}[t]
  \centering
  \anvilfig{fig6-participation.pdf}{fig6\_participation.py}
  \caption{geode-fem stiffness-participation spectrum, with Palace's
  mode locations from the \texttt{Target}~$= 2$\,GHz re-run overlaid
  (its two modes above $33$\,GHz are outside the plotted range). The
  qubit mode appears at $3.45$\,GHz in both solvers (geode-fem
  $p = 0.9942$; Palace port EPR $0.99424$). The $17.49$\,GHz mode has
  $p = 1.000$ in geode-fem but a Palace port EPR of
  $2.49\times10^{-8}$: its port energy carries no port voltage
  (Section~\ref{sec:port-artifact}). Palace's port EPR is not plotted on
  geode-fem's participation axis because the two quantities coincide
  only for uniform-voltage port fields (Section~\ref{sec:agreement}).}
  \label{fig:participation}
\end{figure}

\section{The Forward LOM Chain: From Fields to \texorpdfstring{$E_C$}{E\_C}}
\label{sec:lom}

The lumped-oscillator-model (LOM) branch computes the transmon's
Hamiltonian parameters from the \emph{electrostatic} problem on the same
committed mesh --- and it is the chain the remainder of this paper
differentiates. The mesh's metal is split into three node-disjoint
conductors (ground+resonator, feedline, island); a tensor-$\varepsilon$
electrostatic solve yields the Maxwell capacitance matrix by the
energy-reaction form $C_{ij} = \phi_i^{\mathsf{T}} K \phi_j$; the island
self-capacitance is reduced for the floating feedline
($C_\Sigma = C_{ii} - C_{if}^2/C_{ff}$); and the charging energy follows
as $E_C/h = e^2/(2 C_\Sigma h)$, with the transmon spectrum evaluated
charge-basis-exactly after \citet{koch2007}.

Measured on the committed fixture
(\path{benchmarks/transmon_quantum/results.toml}): $C_\Sigma =
136.7$\,fF, hence $E_C = 0.142$\,GHz, $E_J/E_C = 77.6$ (with $E_J$ from
the junction inductance), $\omega_{01} = 3.38$\,GHz, and anharmonicity
$\alpha = -0.158$\,GHz, with the $\alpha \approx -E_C$ sanity holding.
The chain runs end-to-end from committed geometry. These are
coarse-mesh numbers and they are not mesh-converged, as the next
paragraph shows. The linear eigenproblem
resolves the same oscillator at $3.45$\,GHz
(Section~\ref{sec:qubit-mode}); $\omega_{01}$ is its anharmonic
$0 \to 1$ transition, computed from a $C_\Sigma$ that omits the
junction's own $5.5$\,fF, so the two numbers are not expected to
coincide.

\paragraph{Mesh convergence of $C_\Sigma$.}
A convergence study on the committed geometry
(\path{benchmarks/transmon_diffopt/c_sigma_convergence.toml}) refines
the committed mesh uniformly, splitting every tetrahedron into eight,
and re-solves $C_\Sigma$ of the unmodified device, written $X^0$ below.
It uses the scalar-$\varepsilon$ model, the one the shape gradient
differentiates, with linear (P1) elements at $h$, $h/2$ and $h/4$ and
quadratic (P2) elements at $h$ and $h/2$
(Table~\ref{tab:convergence}, Figure~\ref{fig:convergence}a). The $h/4$
mesh has $8{,}532{,}096$ tetrahedra. A sparse direct factorization does
not fit safely at that size on the 96\,GiB host, so $h/4$ and P2 at
$h/2$ are solved by conjugate gradients with an algebraic-multigrid
preconditioner; where both solvers run, they agree to a relative
$10^{-9}$ or better.

\begin{table}[t]
  \centering
  \caption{Mesh convergence of $C_\Sigma$ (fF) for the unmodified device
  under uniform refinement of the committed mesh. Every entry is a
  conforming energy minimum and therefore an upper bound on the
  converged value for the committed polygonal geometry. No P2 solve was
  run at $h/4$, and the tensor-$\varepsilon$ model has no P2 path.
  Source:
  \texttt{benchmarks/transmon\_diffopt/c\_sigma\_convergence.toml}.}
  \label{tab:convergence}
  \small
  \begin{tabular}{@{}lrrrr@{}}
    \toprule
    Level & Tetrahedra & P1, scalar $\varepsilon$ & P2, scalar $\varepsilon$ & P1, tensor $\varepsilon$ \\
    \midrule
    $h$   & $133{,}314$     & $137.706752$ & $90.044141$ & $136.684731$ \\
    $h/2$ & $1{,}066{,}512$ & $101.359397$ & $83.015285$ & $102.375301$ \\
    $h/4$ & $8{,}532{,}096$ & $87.753422$  & n/a         & $89.493342$ \\
    \bottomrule
  \end{tabular}
\end{table}

\begin{figure}[t]
  \centering
  \anvilfig[0.98]{fig7-convergence.pdf}{fig7\_convergence.py}
  \caption{Mesh convergence of $C_\Sigma$. Every plotted point is a
  conforming upper bound on the converged value of its geometry; the
  shaded bands are extrapolated estimates and are not bounds. (a) The
  unmodified device under uniform refinement, P1 and P2, with the
  rigorous upper bound of $83.02$\,fF, the estimate band of
  $74.14$--$79.62$\,fF and the $89.9$\,fF anchor
  (\texttt{c\_sigma\_convergence.toml}). (b) The two optimized designs
  of Section~\ref{sec:refined} over seven held-out adaptive refinement
  rounds with P2 elements. The design targeted at $89.9$\,fF converges
  below $89.02$\,fF; the resolution-corrected design converges below
  $90.08$\,fF, and its estimate band contains the anchor by
  construction (\texttt{refined\_optimization/*/summary.toml}).}
  \label{fig:convergence}
\end{figure}

\emph{Rigorous.} The P1 spaces are nested under the same Dirichlet and
floating constraints, the P2 space on a mesh contains the P1 space on
it, and $C_\Sigma$ is an energy minimum. Every solved value is
therefore an upper bound on the converged $C_\Sigma$ of this geometry.
The smallest, P2 at $h/2$, puts the converged scalar-$\varepsilon$
value of $X^0$ below $83.02$\,fF. The coarse $137.706752$\,fF
overstates it by at least $54.69$\,fF, which is at least $65.8\%$ of the
converged value. There is no rigorous lower bound.

\emph{Estimate, not a bound.} Four extrapolations span
$74.14$--$79.62$\,fF: Richardson extrapolation of the P1 levels at the
observed order and at order one, a fit of $C_\infty + a h + b h^2$
through the three P1 levels, and Richardson extrapolation of the two P2
levels at order one. The observed P1 order is $1.41$. We expect order
one, because the potential near the edge of a zero-thickness sheet
behaves like $r^{1/2}$; that expectation is an inference from theory
and not a measurement. The three levels are therefore pre-asymptotic.
The band is not asymptotic and it is not a bound. The P2 value on the
committed mesh, $90.044141$\,fF, happens to fall near the anchor. It is
not converged either: one refinement lowers it by a further
$7.02$\,fF.

\emph{Tensor permittivity.} With the sapphire tensor, the P1 values in
Table~\ref{tab:convergence} fall in the same way. The converged tensor
$C_\Sigma$ is below $89.50$\,fF, and two Richardson values put it at
$76.61$--$81.75$\,fF, again an estimate. The scalar-minus-tensor
difference changes sign under refinement: it is $+0.75\%$, $-0.99\%$
and $-1.94\%$ of the tensor value at $h$, $h/2$ and $h/4$. The $0.75\%$
scalar-versus-tensor difference is thus a coarse-mesh number and not a
converged model difference.

\emph{Scope.} ``Converged'' means converged for the polygonal surfaces
of the committed mesh, which refinement does not move. It is not a
statement about a regenerated or re-meshed DeviceLayout device.

\emph{Consequences for the chain.} Every $C_\Sigma$, $E_C$, $E_J/E_C$,
$\omega_{01}$ and $\alpha$ quoted at the start of this section is a
coarse-mesh number. The coarse tensor $C_\Sigma$ overstates the
converged one by at least $47.18$\,fF, and because
$E_C \propto 1/C_\Sigma$ the converged $E_C$ is at least $52.7\%$ above
the coarse $0.142$\,GHz. The Koch spectrum has not been re-evaluated at
a converged $C_\Sigma$.

\paragraph{The gap to the design anchor was discretization error.}
The coarse $C_\Sigma = 136.7$\,fF \emph{exceeds} the ${\sim}90$\,fF
design-anchor value back-solved from the blog workflow's expected
$E_C = 0.2156$\,GHz. Earlier versions, following the benchmark as then
committed, diagnosed this as a model-definition difference and not a
truncation artifact, on two checks: grounding the exterior wall moves
$C_\Sigma$ by a relative $6\times10^{-5}$, and the scalar-versus-tensor
sapphire treatment shifts it $0.75\%$ on the coarse mesh. Neither check
tests mesh convergence, and the diagnosis is withdrawn. The convergence
study settles the question. The converged scalar-$\varepsilon$
$C_\Sigma$ of $X^0$ is below $83.02$\,fF, which is at least
$6.88$\,fF, or $7.6\%$, \emph{below} the $89.9$\,fF anchor; the
estimate band puts it $11.4$--$17.6\%$ below. The coarse mesh
overstates $C_\Sigma$ by at least $54.69$\,fF, more than the whole
$47.8$\,fF by which the coarse value exceeds the anchor. The gap was
discretization error, and more than all of it.

What remains is a gap of the opposite sign: converged, geode-fem's
electrostatic $C_\Sigma$ is below the anchor. As an inference only, the
junction's $5.5$\,fF self-capacitance, which is outside this
$C_\Sigma$, closes part of it: added to the estimate band it gives
$79.64$--$85.12$\,fF, and added to the rigorous bound it gives
$88.52$\,fF, both still below the anchor. The rest is not explained
here; candidates are the geometry representation and the anchor's own
derivation. We report the measured values and the discrepancy; we do
not retrofit to the anchor. Section~\ref{sec:optdemo} uses the anchor
as an optimization target and reports what that does and does not
establish.

\section{Differentiable Sensitivities: The Discrete-Adjoint Layer}
\label{sec:adjoint}

\subsection{The \texorpdfstring{$2\times 2$}{2x2} sensitivity matrix}
\label{sec:adjoint-matrix}

For a solve $A(\varepsilon, X)\,x = b(X)$ and a scalar observable
$g(x)$, the discrete adjoint $A^{\mathsf{T}}\lambda = \partial g/\partial
x$ turns the parameter gradient into explicit contractions:
\begin{equation}
  \frac{dg}{d\varepsilon_k}
    = -\lambda^{\mathsf{T}}\frac{\partial A}{\partial \varepsilon_k}x,
  \qquad
  \frac{dg}{dX}
    = \lambda^{\mathsf{T}}\frac{\partial b}{\partial X}
      - \lambda^{\mathsf{T}}\frac{\partial A}{\partial X}x
  \label{eq:adjoint}
\end{equation}
(real case; the complex driven case takes $2\,\mathrm{Re}[\cdot]$ of the
same contractions with the \emph{transpose}, not conjugate-transpose,
solve --- the pencil is complex-symmetric, so the adjoint solve reuses
the forward LU factorization). Everything on the right-hand side is
assembly-level and therefore differentiable on the tensor substrate: the
element Jacobians $\partial A/\partial X$ and $\partial b/\partial X$ are
\emph{exact}, computed by forward-mode dual numbers through the same
closed-form element kernels the solver assembles --- no finite-difference
truncation enters the gradient. Each gradient costs one forward plus one
adjoint solve sharing a single factorization (asserted programmatically:
\texttt{n\_factorizations} $= 1$).

Four adjoint modules are merged and validated, spanning material and
geometry sensitivities on both operators
(Table~\ref{tab:adjoint-matrix}). Each carries a committed test that the
analytic gradient match a central finite difference \emph{of the full
independent pipeline} (perturb the parameter, re-assemble, re-solve
through the public solver path, recompute $g$) --- and, load-bearingly,
mutation tests proving the tolerances bite: the classic complex-adjoint
conjugation error ($A^{\mathsf{H}}$ for $A^{\mathsf{T}}$) is rejected by
the FD check, as is a dropped term. The driven-shape case surfaced a
physics subtlety worth recording: the current-source right-hand side
$b = \mathrm{i}\omega\mu_0\!\int\! N\cdot J$ is geometry-dependent, so
the shape gradient carries a $\partial b/\partial X$ term absent from the
material case --- omitting it fails the finite-difference cross-check
outright.

\begin{table}[t]
  \centering
  \caption{The merged $2\times 2$ discrete-adjoint sensitivity matrix.
  Each cell: one forward + one adjoint solve sharing a single
  factorization; element Jacobians exact (forward-mode duals through the
  production kernels); validated against central finite differences of
  the full pipeline, with mutation tests. ``FD rel-err'' is the observed
  agreement of the analytic gradient with the FD reference in the merged
  test suite (test acceptance bounds are looser; see the module sources).
  Modules in \texttt{crates/geode-core/src/}.}
  \label{tab:adjoint-matrix}
  \footnotesize
  \setlength{\tabcolsep}{4pt}
  \begin{tabular}{@{}llll@{}}
    \toprule
    Operator & Parameter & Module (issue / PR) & FD rel-err \\
    \midrule
    scalar electrostatic (SPD) & material $\partial/\partial\varepsilon$ &
      \texttt{adjoint} (\#570 / \#573) & ${\sim}3\times10^{-8}$ \\
    scalar electrostatic (P1) & geometry $\partial/\partial X$ &
      \texttt{shape} (\#571 / \#575) & ${\sim}1\times10^{-9}$ \\
    driven $H(\mathrm{curl})$ & material $\partial/\partial\varepsilon$ &
      \texttt{driven::adjoint} (\#576 / \#579) & $2.3\times10^{-5}$ (worst region) \\
    driven $H(\mathrm{curl})$ & geometry $\partial/\partial X$ &
      \texttt{driven::shape} (\#577 / \#581) & $2.2\times10^{-9}$, $1.2\times10^{-8}$ (two maps) \\
    \bottomrule
  \end{tabular}
\end{table}

The matrix is deliberately $2\times 2$: material and geometry are the two
design-parameter classes, and the scalar and driven operators are the
electrostatic (LOM) and RF (S-parameter/EPR-workhorse) branches of the
solver. The coarse-mesh LOM tiers of Section~\ref{sec:optdemo} exercise
the scalar column end-to-end, and the refined optimization of
Section~\ref{sec:refined} uses the same stationarity identity in a
native-f64 P2 implementation; the driven column is merged and validated at
the operator level and composes with the driven solve that geode-fem's
GPU path accelerates --- RF shape optimization is the natural follow-on,
not exercised here.

\subsection{The \texorpdfstring{capacitance$\to E_C$}{capacitance-to-E\_C} chain, and a vanishing adjoint}
\label{sec:chain}

Composing the shape adjoint with the LOM chain of Section~\ref{sec:lom}
gives $\partial(C, E_C, \alpha)/\partial(\text{geometry})$ (merged as
\texttt{shape::capacitance\_shape\_gradient}, \#583 / PR \#586). The
composition carries an elegant simplification. The capacitance observable
is the stored field energy, $C = 2W = \phi^{\mathsf{T}}K(X)\phi$, and $W$
is \emph{variationally stationary} in the solved potential $\phi$: at the
forward solution the free rows of the cotangent $K\phi$ vanish
(equilibrium), so the adjoint solution $\lambda \approx 0$ and the entire
shape derivative collapses to the explicit-geometry term
\begin{equation}
  \frac{dC}{dX}
  \;=\; \phi^{\mathsf{T}}\,\frac{\partial K}{\partial X}\,\phi
  \qquad\text{(at the solution)},
  \label{eq:hf}
\end{equation}
a Hellmann--Feynman-like result: no adjoint back-solve is needed for the
capacitance itself. The implementation keeps the general adjoint form
(the mutation test measures the implicit part at $<10^{-8}$ of the
total, confirming the collapse numerically rather than assuming it), and
chains through $\partial E_C/\partial C_\Sigma = -e^2/(2C_\Sigma^2 h)$
exactly. Validation is double: against the analytic parallel-plate law
($\partial C/\partial d = -\varepsilon_0\varepsilon_r$ per unit area
under a gap scaling; $\partial C/\partial A = +2\varepsilon_0
\varepsilon_r$ under an area scaling --- opposite signs, both matched)
and against central finite differences of the full pipeline, with
observed relative errors at the $10^{-9}$ level.

\section{Gradient-Based Optimization: A Loop Proof, a Coarse-Mesh
Detour, and a Refinement-Validated Design}
\label{sec:optdemo}

Four tiers, labeled by what each one establishes: a closed-loop proof
on an analytic fixture (Section~\ref{sec:optdemo-plate});
single-parameter runs on the coarse committed device mesh, which
validate the gradient and are otherwise historical
(Section~\ref{sec:optdemo-pad}); a three-parameter optimizer on that
coarse mesh, kept as a cautionary precursor
(Section~\ref{sec:multiparam}); and the real-device result, a
three-parameter optimization on an adaptively refined mesh with
quadratic elements, followed by a held-out refinement series that
bounds and estimates what the design is worth when converged
(Section~\ref{sec:refined}). The first two tiers use damped Newton on
$E_C(\theta) - E_C^{\mathrm{target}}$ (equivalently
$C(\theta)$) with the analytic $dE_C/d\theta$ from
Section~\ref{sec:chain}; each Newton step costs one forward plus one
adjoint solve on a single factorization.

One distinction governs the device tiers. The convergence study of
Section~\ref{sec:lom} shows that the coarse committed mesh overstates
$C_\Sigma$ by at least $54.69$\,fF, and that the converged device lies
\emph{below} the $89.9$\,fF anchor, not above it. The coarse-mesh tiers
were run before that was known. Their gradient validations and
convergence statements are exact statements about the coarse discrete
problem: the finite-difference checks test the derivative of the
discrete $C_\Sigma$, and the fresh solves test self-consistency on one
discretization. None of them is evidence about the converged device,
and the direction they moved in, lowering $C_\Sigma$ toward the anchor,
is the wrong one for it. Only Section~\ref{sec:refined} measures what
happens to a design under refinement. The framing claim is
capability and step-count, not wall clock: a derivative-free loop
(the parameter sweeps of current practice) spends $N_{\mathrm{params}}$
extra forward solves per step to estimate what the adjoint delivers
exactly for one. No timing comparison against HFSS/Palace-based
optimizers is measured or claimed.

\subsection{Loop proof: Newton to a target \texorpdfstring{$E_C$}{E\_C} on a parallel-plate
fixture}
\label{sec:optdemo-plate}

On a parallel-plate capacitor fixture ($\varepsilon_r = 3$; 343 nodes,
1{,}296 tets; $\theta$ scales the gap, so $d = L(1+\theta)$), Newton
iteration drives the design from $C_0 = 120.0$\,fF ($E_C =
0.1614$\,GHz) to the design-anchor target $E_C = 0.2156$\,GHz. Full
Newton converges in \emph{one} step --- and we state the reason rather
than dress it up: $C \propto 1/(1+\theta)$ makes $E_C(\theta)$
\emph{affine} in $\theta$, so one-step convergence is expected, and this
tier proves the loop end-to-end (gradient, chain rule, update,
convergence test), not a hard optimization. The converged $\theta =
0.335657909$ matches the closed-form optimum to $3\times10^{-15}$. The
honesty check that matters: a \emph{fresh, independent} forward solve at
the converged geometry --- assemble, solve, extract, no adjoint machinery
--- reproduces the target $E_C$ to a relative error of
$1.4\times10^{-15}$. A damped run ($\alpha = 0.5$) converges in 13 steps
and supplies the descent illustration
(Figure~\ref{fig:diffopt}a). At the hit target, the Koch-exact spectrum
gives $E_J/E_C = 51.0$, $\omega_{01} = 4.128$\,GHz, $\alpha =
-0.247$\,GHz --- the ${\sim}4$\,GHz qubit regime the design anchor
implies. All trajectory values are committed in
\path{benchmarks/transmon_diffopt/results.toml}.

\begin{figure}[t]
  \centering
  \anvilfig[0.98]{fig5-diffopt.pdf}{fig5\_diffopt.py}
  \caption{Gradient-based optimization to a target $E_C$ via the
  discrete-adjoint shape gradient. (a) Parallel-plate loop proof:
  $|E_C - E_C^{\mathrm{target}}|$ versus Newton iteration; full Newton
  converges in one step (the parametrization is affine --- expected, and
  disclosed), the damped $\alpha = 0.5$ run supplies the 13-step descent
  curve; the converged design is confirmed by an independent fresh
  forward solve at relative error $1.4\times10^{-15}$
  (\texttt{benchmarks/transmon\_diffopt/results.toml}). (b) Historical,
  on the coarse 133k-tet DeviceLayout mesh under the fixed-topology island-scale
  parameter $\theta$: a within-budget demonstration target
  ($C_\Sigma = 137.0$\,fF) converges in two genuine (non-affine) Newton
  steps, fresh-solve confirmed at $5.6\times10^{-6}$; the bounded run
  toward the $89.9$\,fF design anchor stalls honestly at the mesh
  validity budget $\theta_{\mathrm{safe}} = -0.0073$ (first tetrahedron
  inversion at $\theta = -0.0097$), $33\times$ short of the
  $\theta \approx -0.241$ the anchor requires on that mesh
  (\texttt{benchmarks/transmon\_diffopt/pad\_results.toml}). The
  shortfall is an artifact of the unconverged mesh: converged, the
  device is already below the anchor (Section~\ref{sec:lom}), and
  Section~\ref{sec:refined} reports the refined optimization.}
  \label{fig:diffopt}
\end{figure}

\subsection{The coarse device mesh: validated gradients and a shortfall that was a mesh artifact}
\label{sec:optdemo-pad}

The second tier runs on the coarse committed 133k-tet SingleTransmon
mesh, with a fixed-topology shape parameter $\theta$ that scales the
island conductor in-plane about its centroid (101 island nodes move;
the mesh connectivity is fixed). Two results stand as statements about
the coarse discrete problem, and a third is historical
(\path{benchmarks/transmon_diffopt/pad_results.toml}):

\paragraph{The gradient is real on the device mesh.} At $\theta = 0$ the
adjoint gradient $\partial C_\Sigma/\partial\theta = 198.198$\,fF is
validated against a central finite difference of the \emph{full
independent pipeline} (move the island nodes, re-assemble, re-solve,
re-extract): relative error $1.15\times10^{-4}$ at $h = 10^{-4}$, with a
clean $O(h^2)$ truncation sweep ($1.2\times10^{-2} \to 1.0\times10^{-3}
\to 1.15\times10^{-4} \to 1.0\times10^{-5}$ as $h$ falls
$10^{-3} \to 3\times10^{-5}$). This was the first real-geometry
validation: the same machinery that passed at $10^{-9}$-level on
analytic fixtures holds up on production geometry. It is a statement
about the coarse discrete problem, and the refined-mesh checks of
Section~\ref{sec:refined} supersede it as the load-bearing one.

\paragraph{The loop genuinely converges on the device mesh.} A
demonstration target within the mesh's validity budget ($C_\Sigma =
137.0$\,fF) converges in two Newton steps --- genuinely non-affine now
(the gradient changes by ${\sim}21\%$ along the path) --- and the
converged design
is confirmed by a fresh, independent multi-conductor extraction at
relative error $5.6\times10^{-6}$ (Figure~\ref{fig:diffopt}b). The
entire coarse-mesh run (validation sweep plus both optimization runs)
takes $10.7$\,s.

\paragraph{Historical: the anchor ``shortfall'' was an artifact of an
unconverged mesh.}
On the coarse mesh the unmodified device has $C_\Sigma = 137.7$\,fF, so
the anchor of Section~\ref{sec:lom} ($89.9$\,fF; $E_C = 0.2155$\,GHz ---
the last digit differs from the $0.2156$\,GHz quoted elsewhere because
this run targets the rounded $89.9$\,fF rather than the back-solved
$89.843$\,fF; both values are artifact-faithful) looked like a
reduction of $47.8$\,fF. Two single-parameter attempts at it stalled on
mesh validity, and earlier versions reported them as the paper's honest
negative.

\emph{Rigid island scale.} The $\theta = 0$ gradient puts the anchor at
$\theta \approx -0.241$. The fixed-topology map inverts its first
tetrahedron at $\theta = -0.0097$, and the run's conservative budget
(minimum tet-volume ratio $0.25$) caps
$\theta_{\mathrm{safe}} = -0.0073$, where the bounded Newton run stalls
at $C_\Sigma = 136.5$\,fF. By a local linear estimate that is
$33\times$ short.

\emph{Harmonic mesh motion.} Extending the prescribed island velocity
into the volume with a P1 Laplace solve spreads the deformation. The
field $D$ is computed once ($0.34$\,s; the whole run takes $7.8$\,s),
giving $X(\theta)=X_0+\theta D$. At the same volume-ratio threshold the
safe negative step grows from $-0.007258$ to $-0.064609$, an
$8.90\times$ extension. The initial derivative is $209.087$\,fF per
unit $\theta$, with a full-pipeline central-FD relative error of
$1.054\times10^{-4}$. Bounded Newton reaches $128.792998$\,fF,
confirmed by a fresh extraction, and stalls at the safe boundary,
leaving $38.892998$\,fF to the anchor. The artifact's initial-slope
extrapolation gives $\theta_{\mathrm{anchor}}=-0.2286$, or $3.54\times$
the harmonic budget. The sensitivity falls along the measured path,
from $209.1$ to $114.5$\,fF per unit $\theta$, so \emph{if the
sensitivity does not recover beyond the measured endpoint} the total
is at least $6.25\times$ the safe budget (computed $6.256959$,
floored). That is a conditional extrapolation, not a measured bound.

These are correct statements about the coarse discrete problem, and the
mesh-validity limit of a fixed-topology morph is real. The shortfall
was not. It was an artifact of an unconverged mesh: converged, the
unmodified device is already below the anchor (Section~\ref{sec:lom}),
so the $47.8$\,fF reduction these runs attempted is not one the device
needs, and the $33\times$ and $6.25\times$ factors describe no physical
design problem. The refined optimization of Section~\ref{sec:refined}
moves the other way, by about $11.5$\,fF, and stays well inside the
same volume-ratio floor. Both coarse runs use scalar trace-averaged
sapphire.

\subsection{A cautionary precursor: three parameters on the coarse mesh}
\label{sec:multiparam}

Before the convergence study, a three-parameter optimizer was run on
the coarse mesh
(\path{benchmarks/transmon_diffopt/multiparam_gradient.toml},
\path{benchmarks/transmon_diffopt/multiparam_results.toml}). Three
harmonic morph fields change the island length, the pad width and the
cutout gap; each is exactly zero on the junction nodes and the island
leads, so no parameter moves the junction region. Their gradients agree
with full-pipeline finite differences to relative errors of
$7.9\times10^{-7}$, $3.4\times10^{-10}$ and $1.1\times10^{-8}$. A
bounded Gauss--Newton optimizer with per-step volume-stiffened
re-morphing then drove the coarse-mesh scalar-$\varepsilon$ $C_\Sigma$
from $137.706752$ down to $89.900001$\,fF in 18 steps, within a
$10$\,kHz tolerance on $E_C$ and with fresh-solve agreement of
$1.39\times10^{-14}$. The stiffening exponent was selected post hoc,
and a neighboring exponent reaches a different geometry in 14 steps, so
the design is not unique.

Every one of those checks passed, and as a device design the result is
wrong: \emph{the optimizer moved the wrong way}. It lowered $C_\Sigma$
by $47.8$\,fF, whereas the converged device starts below the anchor and
needs an increase. Under refinement the capacitance of that design
falls to $66.609133$\,fF (P1, $h/2$), $56.298812$\,fF (P1, $h/4$) and
$52.368047$\,fF (P2, $h/2$)
(\path{benchmarks/transmon_diffopt/c_sigma_convergence.toml}). Its
converged value is therefore below $52.37$\,fF, at least $41.7\%$ below
the anchor, with an extrapolated estimate of $45.98$--$48.11$\,fF. The
lesson is methodological. Finite-difference checks and fresh solves
certify consistency on one discretization and say nothing about
discretization error. What survives from this tier is its machinery:
the junction-pinned morph fields, which the refined optimization
reuses.

\subsection{The real device: a refinement-validated optimization}
\label{sec:refined}

This is the paper's real-device design result. Its scope is the one the
committed artifacts state
(\path{benchmarks/transmon_diffopt/refined_optimization/}): the
scalar-$\varepsilon$ electrostatic island capacitance with a neutral
floating feedline, on the polygonal geometry of the committed mesh,
with the junction's $5.5$\,fF self-capacitance excluded. It is not an
optimization of the full-wave eigenfrequency or of the nonlinear
Hamiltonian.

\paragraph{Method.}
Ten cycles of adaptive refinement, marked by a flux-jump indicator and
solved with P1 elements, turn the committed mesh into a
$448{,}279$-tetrahedron mesh. The optimizer then works on that mesh
with P2 fields, solved by conjugate gradients with an
algebraic-multigrid preconditioner. The geometry map is
$X(\theta) = X^0 + \sum_i \theta_i D_i$, with the three junction-pinned
harmonic fields of Section~\ref{sec:multiparam} (island length, pad
width, cutout gap) computed on the coarse mesh and carried to each
refined mesh by exact nodal prolongation. Connectivity, conductor
membership and material values stay fixed, and refinement happens
outside the differentiable map. The gradient is an equilibrium-energy
shape derivative. For a unit island excitation the floating feedline
sits at $-C_{fi}/C_{ff}$, the energy of the effective field
$\phi_i - (C_{fi}/C_{ff})\,\phi_f$ is $C_\Sigma$, and stationarity
removes the implicit field derivatives, as in Eq.~\eqref{eq:hf}. For a
morph velocity $v$ the P2 derivative integrates
$\nabla\!\cdot v\,\lvert\nabla\phi\rvert^2
 - 2\,\nabla\phi^{\mathsf T}\,\nabla v\,\nabla\phi$
with exact four-point quadrature on each affine tetrahedron. This numerical path is native f64 and does not run
on the Burn tape. Accepted steps must reduce the freshly solved
capacitance residual and keep the minimum element-volume ratio above
$0.25$.

\paragraph{Primary run.}
On the adapted mesh the unmodified device has $C_\Sigma =
78.387918$\,fF, below the anchor. The optimizer therefore
\emph{increases} $C_\Sigma$, and reaches $89.900518$\,fF in three
accepted steps
(\path{benchmarks/transmon_diffopt/refined_optimization/p2_cycles10/summary.toml}).
The final geometry has island length $627.323864\,\mu$m and pad width
$29.087261\,\mu$m, from $625.5$ and $24\,\mu$m. The minimum
volume ratio over the accepted steps is $0.61413$. At the final design
all three sensitivities are checked against central differences of two
fully re-solved perturbed geometries; the worst relative discrepancy is
$3.87\times10^{-7}$. The target is reached on that discretization, and
only there.

\paragraph{Held-out refinement.}
The final geometry is then refined for seven further rounds with the
same marking rule, with a fresh P2 solve each round. None of it feeds
the optimizer. Table~\ref{tab:heldout} and
Figure~\ref{fig:convergence}b give the series. The capacitance falls
monotonically from $89.900518$\,fF at $448{,}279$ tetrahedra to
$89.014998$\,fF at $3{,}289{,}543$, and it is still falling, by
$0.05$\,fF in the last round.

\begin{table}[t]
  \centering
  \caption{Held-out refinement of the optimized designs. Round~0 is the
  optimization mesh; no later round feeds the optimizer. ``Design'' is
  the primary run, targeted at $89.9$\,fF; $X^0$ is the unmodified
  device on the same meshes, and the increment is their difference.
  ``Corrected'' is the resolution-corrected design on its own meshes.
  All capacitances are in fF with P2 elements, and each is an upper
  bound on the converged value of its geometry. Source: the two
  \texttt{summary.toml} files under
  \texttt{benchmarks/transmon\_diffopt/refined\_optimization/}.}
  \label{tab:heldout}
  \footnotesize
  \setlength{\tabcolsep}{4pt}
  \begin{tabular}{@{}rrrrrrr@{}}
    \toprule
    Round & Tetrahedra & Design & $X^0$ & Increment & Tetrahedra & Corrected \\
    \midrule
    0 & $448{,}279$ & $89.900518$ & $78.387918$ & $11.512600$ & $448{,}279$ & $90.980104$ \\
    1 & $528{,}324$ & $89.662755$ & $78.177697$ & $11.485058$ & $528{,}677$ & $90.736602$ \\
    2 & $684{,}646$ & $89.470366$ & $78.011928$ & $11.458438$ & $684{,}107$ & $90.538648$ \\
    3 & $900{,}226$ & $89.324261$ & $77.882041$ & $11.442221$ & $905{,}316$ & $90.391037$ \\
    4 & $1{,}239{,}090$ & $89.216843$ & $77.786774$ & $11.430069$ & $1{,}247{,}811$ & $90.281356$ \\
    5 & $1{,}693{,}481$ & $89.128359$ & $77.710173$ & $11.418186$ & $1{,}708{,}342$ & $90.191606$ \\
    6 & $2{,}369{,}256$ & $89.065595$ & $77.653890$ & $11.411705$ & $2{,}398{,}139$ & $90.127899$ \\
    7 & $3{,}289{,}543$ & $89.014998$ & $77.609945$ & $11.405054$ & $3{,}331{,}618$ & $90.076691$ \\
    \bottomrule
  \end{tabular}
\end{table}

\emph{Rigorous.} The held-out spaces are nested and conforming and the
capacitance is an energy minimum, so every value is an upper bound on
the converged capacitance of this polygonal geometry. The converged
value is below $89.02$\,fF, which is at least $0.98\%$ below the
anchor. There is no rigorous lower bound.

\emph{Estimate, not a bound.} The decrement ratios show no trend toward
zero. A geometric tail at the largest and the smallest observed ratio
gives $88.778547$ to $88.891530$\,fF, and least-squares fits of
$C_\infty + a N^{-\alpha}$ with a free rate on three windows of rounds
give $88.900718$, $88.863790$ and $88.892726$\,fF. The outward-rounded
union is $88.77$--$88.91$\,fF, about $88.8$--$88.9$\,fF. On that
estimate the optimized $89.900518$\,fF overstates the converged value
by $1.11$--$1.28\%$. The estimate depends on the assumed rate: the same
fit at slower assumed rates gives $88.336686$\,fF ($\alpha = 1/3$) and
$88.618200$\,fF ($\alpha = 1/2$), below the band. That is why the band
is an estimate and not a lower bound.

\emph{No $1\%$ claim.} The rigorous statement, at least $0.98\%$ below
the anchor, does not decide whether the design is within $1\%$ of it.
The estimate puts it $1.10$--$1.26\%$ below. We do not claim the $1\%$
threshold is met, and we infer from the extrapolation that it is not.

\paragraph{The design increment is stable.}
On the same meshes the unmodified device goes from $78.387918$ to
$77.609945$\,fF, so its converged value is below $77.61$\,fF, at least
$13.6\%$ below the anchor, and is estimated at $77.42$--$77.52$\,fF.
That is inside the $74.14$--$79.62$\,fF band of the uniform study and
tighter than its $83.02$\,fF bound (Section~\ref{sec:lom}). The
increment the optimizer added moves much less than either capacitance:
it is $11.512600$\,fF on the optimization mesh and $11.405054$\,fF
after seven rounds, a change of less than $1\%$. Most of the
discretization error is common to the two geometries on these meshes,
which are adapted to the final design and not to $X^0$. That is an
observation about this series, not a bound.

\paragraph{A resolution-corrected design, and what it does not show.}
A second run repeats the optimization on the same discretization with
the target offset by the measured bias: the anchor times the
optimization-mesh value over the geometric-tail midpoint,
$89.9 \times 89.900518 / 88.835038 = 90.978253$\,fF. The first run's
artifact records that target, which was fixed before the second run
(\path{benchmarks/transmon_diffopt/refined_optimization/p2_cycles10_resolution_corrected/summary.toml}).
The optimizer reaches $90.980104$\,fF in three accepted steps (island
length $627.452356\,\mu$m, pad width $29.447205\,\mu$m, minimum volume
ratio $0.57790$, worst finite-difference discrepancy
$3.94\times10^{-7}$). Its held-out series runs from $90.980104$ to
$90.076691$\,fF at $3{,}331{,}618$ tetrahedra.

\emph{Rigorous.} The converged capacitance of this geometry is below
$90.08$\,fF, at most $0.20\%$ above the anchor. The last solved value is
still above the anchor, so nothing rigorous places the converged value
near it from below.

\emph{Estimate, and a caveat that must be stated plainly.} The same two
methods give a band of $89.84$--$89.96$\,fF, within $0.07\%$ of the
anchor on either side. \textbf{This band contains the anchor by
construction.} The corrected series is the first series rescaled:
round by round, the ratio of the two is constant to within
$8.1\times10^{-5}$ (our computation from the two committed series). An
extrapolation rule applied to both therefore returns, to about that
accuracy, the first run's extrapolated value times that ratio, and the
target was set to the
anchor times the inverse of the first run's extrapolated bias. The
width of the band is the spread between extrapolation methods. It is
not an uncertainty. What the second run does establish is the upper
bound above, and that the relative discretization bias carries over
between two nearby designs to about $10^{-4}$. At slower assumed rates
the same series extrapolates to $89.394848$\,fF ($\alpha = 1/3$) and
$89.678525$\,fF ($\alpha = 1/2$), which is $0.57\%$ and $0.25\%$ below
the anchor. Only the upper side is proven. A two-sided bound needs a
guaranteed lower bound, for instance from a complementary-energy
formulation of the same capacitance, or a substantially finer or
higher-order solve. Neither exists here; it is future work (issue
\#1062).

\paragraph{Cross-solver check: the same discretization.}
Palace, run with second-order elements on the exported adapted mesh of
the primary design, gives $89.900294$\,fF against geode-fem's
$89.900518$\,fF, a relative difference of $2.5\times10^{-6}$. On the
resolution-corrected design it gives $90.979877$ against
$90.980104$\,fF, the same relative difference. Palace adds $3{,}777$
elements by local bisection at the interior boundaries, so the two
codes share the polygonal geometry and coefficients, not one mesh. This shows that two solvers agree on the same discretization. It
is not evidence about the discretization error, which the held-out
series measures. A third-order Palace run was killed (signal~9) during
its error estimation and produced no complete capacitance matrix; its
log is retained and it is not evidence. Two adaptive Palace steps left
the island capacitance unchanged while the feedline self-capacitance
moved by $1.90\%$, which suggests the marking concentrated on the
feedline; we do not count them as a convergence check either.

\paragraph{What is not claimed.}
No continuum value is certified for either design. No scalar-to-tensor
equivalence, layout regeneration or re-meshing, full-wave optimization,
experimental validation, global optimum or speedup against Palace is
claimed. The geometry has zero-thickness sheet conductors and a natural
zero-charge outer boundary, and ``converged'' means converged for the
polygonal surfaces of the committed mesh. The choice of P2 elements and
ten refinement cycles was made after exploratory P1 runs, which are
retained in the artifact directory; it was not preregistered.

\paragraph{Roadmap.}
Two extensions are planned and not built. The first is the two-sided
bound above. The second is a second design target, the island--claw
coupling ratio $\beta = C_g/C_\Sigma$ (issue \#1037): the coarse-mesh
optimizer interface accepts several targets and is unit-tested with two
targets and three parameters, but no two-target experiment on the
device exists.

\section{Roadmap: Differentiating the Eigenmode Branch}
\label{sec:roadmap}

The LOM branch above differentiates an \emph{electrostatic} solve. The
transmon's dynamical figures-of-merit --- resonator frequency,
participation-ratio and EPR self/cross-Kerr \citep{minev2021,nigg2012}
--- live in the \emph{eigenmode} problem, and that branch is \textbf{not
differentiable yet} in geode-fem. The paper commits to the explicit
split:

\begin{table}[h]
  \centering
  \small
  \begin{tabular}{@{}lll@{}}
    \toprule
    Branch & Observables & Differentiable? \\
    \midrule
    LOM (electrostatic) & $E_C$, $\alpha \approx -E_C$, coupling $C$ &
      \textbf{yes} (Sections~\ref{sec:adjoint}--\ref{sec:optdemo}) \\
    eigenmode / EPR & resonator $f$, participation, Kerr &
      \textbf{not yet} (this section) \\
    \bottomrule
  \end{tabular}
\end{table}

Differentiating an eigenpair needs the adjoint-eigenpair
(Hellmann--Feynman / Nelson) formulas
\citep{nelson1976simplified}, which require solving the interior
eigenproblem robustly at the physical shift --- and geode-fem's
matrix-free interior eigensolve is blocked at exactly that shift. The
blocker is measured and committed
(\path{benchmarks/transmon_bench_cpu/sigma4p5_deepshift_characterization.md}):
at $\sigma = 4.5$\,GHz on the 133k-DOF fixture, a \emph{single} inner
AMS-preconditioned MINRES solve (outer Lanczos step 0) never converges
--- it plateaus on a nearly flat residual tail below ${\sim}10^{-5}$ and
was still running at 14{,}300 inner iterations ($\lVert r\rVert \approx
2.6\times10^{-6}$) when killed. The plateau is
\emph{coarse-solve-invariant} --- an exact coarse solve stalls the same
way --- so the limiter is the SPD-proxy preconditioner $K + |\sigma|M$
itself, not the coarse solve. Nor is there a single inner tolerance that
is both cheap and correct: at an inner tolerance of $10^{-2}$ the entire
eigensolve completes in ${\sim}56$\,s but returns the non-physical
$\lambda \approx 0$ gradient near-kernel cluster ($0.42$--$0.64$\,GHz),
not the physical band. The named path --- not built, cited as method,
not result --- is a preconditioned Helmholtz-projected Jacobi--Davidson
iteration, applying the gradient-kernel projection inside the correction
equation, for which \citet{liang2026adaptive} provide the current
adaptive-multilevel treatment of exactly this singular Maxwell
eigenproblem; combined with the Nelson/Hellmann--Feynman eigenpair
derivatives, that is the route to differentiable eigenmode/EPR
quantities. Until it lands, no claim in this paper extends to
derivatives of the eigenmode spectrum.

\section{Performance Context}
\label{sec:perf}

The timings below qualify the cost of the capability. Eigenmode,
driven CPU and driven GPU results are separate experiments; none measures
the optimizer's cost or establishes a general solver ranking.

\subsection{CPU eigenmodes: same box, verified mode content}
\label{sec:perf-cpu}

The earlier timing request, six modes at $4.5$\,GHz, did not establish
equivalent work. In the mode-logged re-run it returns four near-kernel
modes, the $3.45$\,GHz qubit mode and only one of Palace's six modes in
geode-fem, whereas Palace, which searches only above its target, returns
its six. The historical run has
no committed mode log or CPU-time record; its mode content and actual
thread count cannot be recovered. We therefore withdraw the earlier
claim that one geode core beats eight Palace ranks with roughly twelve
times fewer core-seconds. This does not change the separate frequency
agreement credential of Section~\ref{sec:agreement}.

The replacement experiment ran on one otherwise idle EC2
\texttt{m6i.4xlarge} in \texttt{us-east-1} on 2026-10-08, with geode
\texttt{61e8571e} and Palace \texttt{fba6a5b}. One-wide runs were
OS-pinned to one CPU; eight-wide runs used one hardware thread per
physical core. Wall clock covers the whole solver process; core-seconds
are measured user plus system CPU time, not wall time multiplied by
requested width. Table~\ref{tab:cpu} retains all three geode requests
that recover the six reference modes. Two are tuned using the reference
spectrum: requesting thirty modes at $4.5$\,GHz, or moving the shift to
$20$\,GHz. The older port-aware extraction route does not use reference
frequencies to choose its request, but returns two further modes
alongside the six (one near-kernel mode and the qubit mode); matching
the six still uses Palace.
Palace's target is a lower bound, whereas geode's shift is a centre, so
identical numerical requests need not select identical spectra.

Two asymmetries qualify every comparison in Table~\ref{tab:cpu}. First,
the solvers do not solve the same number of unknowns. Palace refines
the input mesh around the metal sheets, and every timed Palace run
assembles $183{,}891$ N\'ed\'elec unknowns against geode-fem's
$156{,}863$ ($133{,}108$ interior), about $17\%$ more
(Section~\ref{sec:geometry}). The table does not correct for this, and
the difference favours geode-fem. Second, the reference set is Palace's
six modes above $4.5$\,GHz. That set includes the $17.49$\,GHz
port-sheet mode and excludes the qubit mode, which Palace cannot return
at this target (Section~\ref{sec:honest-physics}). The ``Extra modes''
column counts returned modes outside that set. In the thirty-mode and
port-aware rows one of the extras is the qubit mode, which the committed
artifact still labels \texttt{spurious\_port}; the others are
near-kernel modes.

\begin{table}[t]
  \centering
  \caption{Mode-verified CPU timings on one m6i.4xlarge: median wall
  seconds / measured core-seconds. All rows contain the six reference
  modes of Table~\ref{tab:agreement}. The solvers start from the same
  input mesh but Palace solves $183{,}891$ unknowns and geode-fem
  $156{,}863$; the timings are not normalized for that.
  Three repeats per cell except Palace one-rank Save=0 (one run).
  The two oracle-tuned requests are marked $\dagger$; the historical
  port-aware route is the pre-port-subspace implementation.
  Palace always requests six modes above $4.5$\,GHz.}
  \label{tab:cpu}
  \small
  \begin{tabular}{@{}lrrr@{}}
    \toprule
    Request & Extra modes & 1 wide & 8 wide \\
    \midrule
    geode 30 at 4.5 GHz$^\dagger$ & 24 & 34.45 / 34.44 & 29.08 / 40.71 \\
    geode 6 at 20 GHz$^\dagger$ & 0 & 33.21 / 33.20 & 27.26 / 38.87 \\
    geode port-aware 8 at 4.5 GHz & 2 & 67.26 / 67.25 & 51.12 / 91.72 \\
    Palace, Save=6 & 0 & 186.84 / 186.82 & 33.35 / 266.66 \\
    Palace, Save=0 & 0 & 113.00 / 112.98 & 23.64 / 188.89 \\
    \bottomrule
  \end{tabular}
\end{table}

\begin{figure}[t]
  \centering
  \anvilfig[0.98]{fig4-cpu-wallclock.pdf}{fig4\_cpu\_wallclock.py}
  \caption{All three mode-recovering geode requests beside Palace with
  and without field output, from Table~\ref{tab:cpu}. Bars show medians;
  whiskers show the observed minimum and maximum. The oracle-tuned
  requests are not a reference-free recipe for an unknown device.
  Save=0 removes Palace's field output, but retains its error estimator.}
  \label{fig:cpu}
\end{figure}

At eight wide Palace without field output takes $23.64$\,s, less than
any of these geode routes ($27.26$--$51.12$\,s). With field output,
Palace takes $33.35$\,s and the two oracle-tuned routes are faster.
At one wide the geode routes take less wall and CPU time, but neither
comparison isolates a solver kernel: geode writes no fields and has no
counterpart to Palace's error estimator, which remains enabled even
with Save=0 (about $40.2$\,s at one rank and $5.0$\,s at eight).
MPI busy-waiting also affects measured Palace core-seconds. Palace also
carries about $17\%$ more unknowns in every cell. No single
route is selected as a headline winner.

The Palace build uses SuperLU-preconditioned GMRES, not AMS:
the pinned source and build configuration resolve \texttt{Default} to
SuperLU. A later, separate m6i session with a tuned Palace build found
only a $1.01$--$1.05\times$ same-box wall-time improvement over its
untuned control; cross-session differences must not be credited to
compiler tuning. Those runs also predate the new port-subspace solver.
The local update in Section~\ref{sec:gauge-arc} cannot be substituted
into this same-box timing table.

At $1{,}157{,}564$ DOFs, a historical geode direct run completed on a
128\,GB box in $565.5$\,s with $92.2$\,GB peak RSS. Palace's earlier
eight-rank run recorded $423.12$\,s and about $33$\,GB aggregate memory.
These are different-host observations, and the geode request lacks
mode-content verification. The DOF count is geode-fem's; we infer that
Palace's internal refinement applies to that mesh as well, and its
unknown count for that run is not quoted here. They establish the direct factorization's
large memory requirement, not a controlled performance crossover or
equivalent physical-mode throughput.

\subsection{CPU driven solve: AMS at scale}
\label{sec:perf-ams}

A separate f64 experiment uses assembled COCG on the lossy driven cube
at $\omega=0.10$, on one dedicated \texttt{r6i.4xlarge}, with geode
\texttt{64d44269}. AMS reduces one-core linear-solve time relative to
Jacobi by $3.4$--$4.2\times$ across $25{,}695$--$897{,}650$ edges.
It reduces iterations from $2969$--$13662$ to $39$--$105$, but does
not eliminate iteration growth. Every geode solve meets an explicitly
recomputed relative-residual tolerance of $10^{-8}$.

\begin{table}[t]
  \centering
  \caption{Driven CPU-f64 linear-solve seconds on one r6i.4xlarge,
  one core/rank, median of two geode and three Palace runs, except
  $897{,}650$ edges (one run each). geode includes operator assembly,
  preconditioner setup, iteration and residual check; Palace sums its
  Setup, Preconditioner and Linear Solve timers, excluding Operator
  Construction. These are not whole-process timings.}
  \label{tab:ams}
  \small
  \begin{tabular}{@{}rrrrr@{}}
    \toprule
    Edges & geode Jacobi & geode AMS & Palace AMS & Palace SuperLU \\
    \midrule
    25,695 & 2.4 & 0.7 & 1.0 & 1.4 \\
    102,024 & 22.2 & 5.3 & 5.0 & 8.9 \\
    462,520 & 227.3 & 56.4 & 26.4 & 85.9 \\
    897,650 & 624.3 & 181.3 & 53.3 & 235.8 \\
    \bottomrule
  \end{tabular}
\end{table}

At $462{,}520$ edges AMS needs about $3.4$\,GB versus Jacobi's
$1.05$\,GB. At $897{,}650$ edges it needs $8.5$\,GB versus $2.0$\,GB,
and its setup alone takes $53.488$\,s. Eight threads reduce setup but
do not accelerate geode's Krylov iteration; at $462{,}520$ edges its
AMS solve takes $50.4$\,s versus Palace AMS's $4.2$\,s on eight ranks.
The one-core geode AMS advantage at the smallest size does not persist
at scale: Palace AMS takes $53.3$\,s against $181.3$\,s at the largest
size. These measurements concern one structured first-order fixture,
assembled CPU operators and low frequencies. They establish neither
GPU performance nor a remedy for the interior eigenproblem.

\subsection{GPU cell: correctness proven, scaling an honest negative}
\label{sec:perf-gpu}

geode-fem's GPU path accelerates the \emph{driven} solve (matrix-free
apply + on-device COCG), not the eigensolve of the credential, and all
GPU execution is CUDA-f32 (Section~\ref{sec:architecture}). On
2026-07-14, geode-fem code executed on a physical GPU for the first time
(EC2 \texttt{g6e.xlarge}, 1$\times$ NVIDIA L40S 46\,GB, driver
595.71.05, CUDA 13.2): the two CUDA-f32 correctness smokes --- the
matrix-free N\'ed\'elec matvec and the on-device COCG --- pass against
the f64 CPU reference within f32-appropriate tolerances ($15.0$ and
$3.3$\,s wall, durations dominated by GPU init and kernel JIT ---
correctness evidence, not performance).

The scaling study (committed as
\path{benchmarks/gpu_driven_scaling/results.toml}; a $\sigma$-lossy
lumped-port cube fixture at fixed frequency, mesh-refined from $1{,}854$
to $25{,}695$ edges, all sixteen cells on the same host) is an honest
negative: \textbf{GPU-f32 matrix-free loses to every CPU configuration
at every measured size} (Table~\ref{tab:gpu}). Computed from the
artifact's unrounded medians, the assembled-CSR COCG on CPU is
$136\times$ faster at $1{,}854$ edges (displayed: $0.032$ versus
$4.39$\,s) and still $44\times$
faster at $25{,}695$ ($1.86$ versus $81.8$\,s); even the direct LU is
$13.5\times$ faster at the top size ($6.04$ versus $81.8$\,s). The
mechanism is kernel-launch domination (${\sim}28$\,ms per iteration at
the top size versus ${\sim}0.63$\,ms for assembled-CSR on CPU) ---
$25.7$k edges is far too small to fill an L40S --- and f32 accuracy
compounds the verdict (relative error degrading to $1.2\times10^{-3}$
with refinement; nondeterministic convergence at the top size). One
directional positive survives, and it is the one the architecture
predicts: against the \emph{same matrix-free algorithm on CPU}, the
GPU's deficit falls monotonically to parity at ${\sim}26$k edges
($81.8$ versus $82.1$\,s). We state the GPU position as trajectory, not
achievement, with tracked levers (kernel fusion; on-device
f64/mixed-precision when cubecl ships it; preconditioning) --- and we
prefer publishing the bound to implying the win. This cell is a cube
fixture family, not the transmon eigenproblem, and its CPU baselines
(4-vCPU host) are not comparable to the \texttt{m6i} numbers of
Table~\ref{tab:cpu}.

\begin{table}[t]
  \centering
  \caption{Driven-solve wall-clock scaling, solve-only medians of $n=3$
  timed repetitions (warm-up excluded), same host (\texttt{g6e.xlarge}:
  4 vCPU + 1$\times$ L40S). One fixture ($\sigma$-lossy lumped-port
  cube) at fixed $\omega$, mesh-refined across the columns. f64
  iterative configurations at tolerance $10^{-8}$; the f32 GPU
  configuration at $10^{-6}$ (its recurrence floor sits above
  $10^{-8}$). The accuracy row is the full-field relative $L_2$ error of
  the GPU-f32 solution against the Direct-f64 reference. At 25{,}695
  edges the GPU-f32 configuration converged in the committed
  single-frequency runs but not in its five-point frequency sweep, and
  run-to-run convergence is nondeterministic at that size (f32 floor).
  Values verbatim from the committed
  \texttt{benchmarks/gpu\_driven\_scaling/results.toml}.}
  \label{tab:gpu}
  \small
  \setlength{\tabcolsep}{5pt}
  \begin{tabular}{@{}lcccc@{}}
    \toprule
    & \multicolumn{4}{c}{Mesh size (N\'ed\'elec edges)} \\
    \cmidrule(l){2-5}
    Configuration & 1{,}854 & 5{,}859 & 13{,}428 & 25{,}695 \\
    \midrule
    Direct faer LU (CPU f64), s & 0.024 & 0.203 & 1.540 & 6.036 \\
    Assembled-CSR COCG (CPU f64), s & 0.032 & 0.198 & 0.709 & 1.865 \\
    Matrix-free COCG (CPU f64), s & 1.652 & 8.885 & 30.234 & 82.056 \\
    Matrix-free COCG (GPU f32), s & 4.388 & 13.351 & 34.381 & 81.764 \\
    \midrule
    GPU-f32 accuracy vs.\ Direct-f64 (rel.\ $L_2$) &
      $1.26\times10^{-4}$ & $2.93\times10^{-4}$ &
      $4.60\times10^{-4}$ & $1.20\times10^{-3}$ \\
    \bottomrule
  \end{tabular}
\end{table}

\section{Reproducibility}
\label{sec:repro}

Every artifact behind every result in this paper is committed and
scripted.\footnote{Repository URL and archival DOI for the geode-fem
source to be added at submission. TODO(operator).}
The exceptions are labelled where they occur: the scratch-probe
measurements and source readings of Section~\ref{sec:honest-physics}
trace to a committed design note and probe source, not to a regenerable
benchmark artifact.

\begin{itemize}
  \item \textbf{Geometry and mesh.} Generated by DeviceLayout.jl
    v1.15.0's \texttt{SingleTransmon} example \citep{devicelayout}; the
    fixture generator is committed, and the MSH~4.1 fixture is pinned by
    SHA-256 (\texttt{5b3ff4c3\ldots{}b33dd}; full hash and provenance in
    the committed \texttt{transmon\_smoke.provenance.txt}).
  \item \textbf{The MSH conversion gotcha.} MFEM's MSH~4.1 reader
    rejects the DeviceLayout multi-entity-block file, so the Palace runs
    consume an MSH~2.2 conversion (\texttt{gmsh -save -format msh2}).
    The conversion is node-preserving (all 22{,}684 nodes unchanged,
    physical groups preserved) and physics-neutral: eigenvalues
    reproduce bit-for-bit across the conversion. Documented in the
    committed \texttt{palace\_config.provenance.txt} because it is
    exactly the kind of silent pipeline step a same-mesh claim must
    disclose. A second such step went undisclosed in earlier versions:
    Palace refines the mesh it reads ($4{,}913$ elements added by local
    bisection, $6{,}014$ duplicated vertices) and solves $183{,}891$
    unknowns, as its committed run logs record
    (\path{reference/fixtures/transmon_palace/results_p1/palace_run_v22.log}).
  \item \textbf{Credential artifacts.} The Palace configuration and
    provenance under \path{reference/fixtures/transmon_palace/};
    geode-fem's benchmark definitions, the agreement numbers, the
    L-doubling tripwire ($17.4901 \to 12.37$\,GHz, committed ratio
    $0.7071$), the qubit-mode measurements ($3.4528$\,GHz,
    $p = 0.9942$; L-doubling ratio $0.7081$; stored under the
    artifact's \texttt{[spurious\_mode]} key), Palace's
    \texttt{Target}~$= 2$\,GHz re-run
    (\texttt{[oracles.palace\_target2]}), and the full
    gauge/projection arc of Section~\ref{sec:gauge-arc}
    (tree--cotree counts and the $1.64\%$ shift; bulk-projection
    divergence ratio $50.2$ and projected norm $1.06\times10^{-4}$;
    port-aware acceptance at $\leq 0.029\%$) in
    \path{benchmarks/transmon_eigen/results.toml}, each pinned by a
    release-tier regression test. Palace's raw outputs
    (\texttt{eig.csv}, per-port participation, run log) are committed;
    a rerun reproduced \texttt{eig.csv} bit-for-bit. The re-run's
    configurations, logs, \texttt{eig.csv}, \texttt{port-EPR.csv} and
    provenance are under
    \path{reference/fixtures/transmon_palace/results_p1_target2/}. The
    source reading of Palace's port operator and search window, the
    port-voltage probe and the planned rank-1 formulation are in
    \path{docs/research/transmon-lumped-port-formulation.md}; its probe
    is scratch source that the build does not compile.
  \item \textbf{Precision of the committed agreement numbers.}
    \path{benchmarks/transmon_eigen/results.toml} stores geode-fem
    frequencies at three decimals while the Palace column carries full
    precision, so the committed per-mode \texttt{rel\_err\_pct} is
    computed against the rounded geode-fem values; at full precision the
    agreement is slightly better (approximately $0.029\%$ worst case).
    We quote the committed-artifact numbers throughout and disclose the
    rounding here. Derived ratios quoted in the text (core-second and
    speedup multipliers) are computed from the displayed table values
    unless marked as committed.
  \item \textbf{LOM chain.} The three-conductor capacitance extraction
    and Koch-exact spectrum on the coarse mesh ($C_\Sigma = 136.7$\,fF,
    $E_C = 0.142$\,GHz, and the boundary and sapphire-model checks of
    Section~\ref{sec:lom}) in
    \path{benchmarks/transmon_quantum/results.toml}; pinned by the
    release test \path{crates/geode-core/tests/transmon_quantum.rs}.
    The mesh-convergence study (P1 to $h/4$, P2 to $h/2$, scalar and
    tensor permittivity, with its bounds, estimates and memory record)
    in \path{benchmarks/transmon_diffopt/c_sigma_convergence.toml},
    written by the example \texttt{transmon\_c\_sigma\_convergence}; it
    ran locally in under five minutes; earlier runs of it peaked at
    $12.8$--$13.8$\,GB of resident memory.
  \item \textbf{Adjoint layer.} The four sensitivity modules with their
    FD-and-mutation test suites:
    \path{crates/geode-core/src/adjoint.rs} (PR \#573);
    \path{crates/geode-core/src/shape.rs} (PRs \#575/\#586; this module
    also carries \texttt{capacitance\_shape\_gradient});
    \path{crates/geode-core/src/driven/adjoint.rs} (PR \#579); and
    \path{crates/geode-core/src/driven/shape.rs} (PR \#581).
  \item \textbf{Optimization demonstrations.} The parallel-plate loop
    proof (full trajectories, closed-form check, fresh-solve
    confirmation) in \path{benchmarks/transmon_diffopt/results.toml}.
    The refined real-device optimization under
    \path{benchmarks/transmon_diffopt/refined_optimization/}: each run
    directory (\texttt{p2\_cycles10/} and
    \texttt{p2\_cycles10\_resolution\_corrected/}) holds a generated
    \texttt{summary.toml} with the optimizer steps, the final
    finite-difference checks, the held-out series and every bound and
    estimate with its rounding direction, plus the Palace
    configurations, logs and outputs; \texttt{validation\_summary.json}
    is a hand-assembled index of the same numbers, and the directory's
    \texttt{README.md} gives the commands (about 20 and 10 minutes
    locally, with peak resident memory up to $29.4$ and $26.0$\,GB).
    The coarse-mesh runs: the single-parameter FD validation and
    bounded anchor attempt in
    \path{benchmarks/transmon_diffopt/pad_results.toml}; harmonic motion
    in \path{benchmarks/transmon_diffopt/harmonic_results.toml}; the
    three-parameter gradients in
    \path{benchmarks/transmon_diffopt/multiparam_gradient.toml}; and the
    bounded optimizer runs in
    \path{benchmarks/transmon_diffopt/multiparam_results.toml}. All are
    consumed by \path{crates/geode-core/tests/transmon_diffopt.rs},
    which pins both \texttt{summary.toml} files and re-derives their
    bounds and estimates.
  \item \textbf{Performance artifacts.} The historical CPU cell in
    \path{benchmarks/transmon_bench_cpu/results.toml}; mode-verified
    timing in
    \path{benchmarks/transmon_bench_cpu/results_like_for_like_m6i.toml};
    tuned-build controls in
    \path{benchmarks/transmon_bench_cpu/results_tuned_palace_m6i.toml};
    local port-subspace evidence in
    \path{benchmarks/transmon_bench_cpu/results_like_for_like_local.toml};
    driven CPU AMS results in
    \path{benchmarks/gpu_driven_scaling/results_ams_cpu_r6i.toml}; the 1.16M-DOF
    completion run log at
    \path{benchmarks/transmon_bench_cpu/geode_runs_1p16M_2026-07-15.log};
    the GPU scaling cell (all sixteen cells, iteration counts,
    recomputed true residuals, DNF handling, same-host provenance) in
    \path{benchmarks/gpu_driven_scaling/results.toml}; the deep-shift
    eigensolve characterization at
    \path{benchmarks/transmon_bench_cpu/sigma4p5_deepshift_characterization.md}.
  \item \textbf{Hardware disclosure.} CPU cell: EC2 \texttt{m6i.4xlarge}
    (8 physical cores / 16 vCPU, 64\,GB), \texttt{us-east-1} for the
    mode-verified re-timing (the historical session was in
    \texttt{us-west-2}); driven AMS: \texttt{r6i.4xlarge},
    \texttt{us-east-1}; the
    1.16M-DOF completion row: a 128\,GB \texttt{r6i}-class box (see the
    committed log). GPU cell: EC2 \texttt{g6e.xlarge} (4 vCPU, 30\,GB;
    1$\times$ NVIDIA L40S 46\,GB; driver 595.71.05, CUDA 13.2). The
    Palace \texttt{Target}~$= 2$\,GHz re-run: a local Apple M3 Ultra,
    Palace built from the in-repository Docker recipe (arm64), eight MPI
    ranks; no timing from it is used.
  \item \textbf{Versions.} geode-fem commit \texttt{3174015} (headline
    credential and historical CPU cells; the adjoint/diffopt layers are merged
    post-\texttt{3174015} and cited by PR number); Palace commit
    \texttt{fba6a5b}; DeviceLayout.jl v1.15.0; Burn/burn-cuda 0.21.
\end{itemize}

\section{Discussion}
\label{sec:discussion}

\paragraph{What the credential does and does not establish.}
In V\&V terms \citep{roache1997quantification,oberkampf2010verification}
the cross-validation is code verification and code-to-code comparison,
not validation against nature: two independent implementations agreeing
to $\leq 0.033\%$ from the identical input mesh, on discrete problems
that differ by Palace's internal refinement, with scaling tripwires on
the port operator and on the qubit mode. It does not certify that either
solver matches a measured device --- that is what
\citet{ye2025electromagnetic}-style experimental benchmarks address ---
but it rules out the large error class living in formulation, assembly,
boundary handling, and eigensolver machinery. It establishes that
geode-fem's eigenmode path reproduces a production solver from the same
input. It does not establish mesh convergence. In particular the
electrostatic observable that this paper differentiates is not
converged on this mesh (Section~\ref{sec:lom}), and the eigenmode
credential says nothing about that.

\paragraph{What the differentiable-design evidence supports.}
The proven capability is bounded and we state its boundary. Proven: the
$2\times 2$ adjoint matrix at operator level (FD-validated with mutation
tests); the capacitance$\to E_C$ chain with its stationarity collapse;
a closed Newton loop confirmed by an independent fresh solve on an
analytic fixture; a three-parameter optimization of the real device on
an adaptively refined P2 discretization, with gradients validated
against fully re-solved finite differences at $3.87\times10^{-7}$ or
better; and a rigorous upper bound on the converged capacitance of each
optimized design (below $89.02$\,fF, and below $90.08$\,fF for the
resolution-corrected one; Section~\ref{sec:refined}).
Estimated, not proven: the converged values themselves (about
$88.8$--$88.9$\,fF and $89.84$--$89.96$\,fF). The second band contains
the anchor by construction.
Not proven, and not claimed: a lower bound on either converged
capacitance, and therefore any certified distance to the anchor from
below; that the first design is within $1\%$ of the anchor;
eigenmode/EPR derivatives (the Section~\ref{sec:roadmap} wall);
two-target or coupling-constrained optimization;
tensor-$\varepsilon$ shape chains; a design regenerated by the layout
toolchain or re-meshed; and any wall-clock superiority
over sweep-based workflows. The capability is also not qubit-specific
--- shape and material sensitivities of electrostatic and driven-EM
observables apply to RF and photonic device design generally; the
transmon is the demonstration vehicle with the clearest documented
unmet need.

\paragraph{Threats to validity.}
A common-mode error --- both solvers wrong the same way --- cannot be
excluded by cross-comparison alone, and this paper now contains an
example. The $17.49$\,GHz mode is an artifact of a port representation
that both solvers share; the two agree on it to $0.001\%$, and earlier
versions took that agreement, with the $f_{LC}$ anchor, as confirmation
of the junction physics. The same versions asserted the absence of a
mode that Palace had not been asked for. Agreement with a reference
bounds implementation error, not shared modeling choices or the
reference's own search window. Every capacitance in this paper is a
conforming upper bound, and no lower bound has been computed. The
estimates of converged values rest on extrapolations whose rates are
not established; at slower assumed rates the resolution-corrected
design lands $0.25$--$0.57\%$ below the anchor. That design's estimate
band contains the anchor by construction. The two refined runs share a
mesh family, a marking rule and an element order, so neither confirms
the other independently, and Palace's agreement on the adapted mesh is
agreement on the same discretization. The choice of P2 elements and ten
refinement cycles was made after exploratory runs and was not
preregistered. No two-target or coupling-constrained optimization is
demonstrated. The
parallel-plate tier's one-step convergence is an affine-map property,
disclosed as such. The real-device chain uses the scalar trace-averaged
sapphire, and the scalar-versus-tensor difference is not converged: it
changes sign under refinement (Section~\ref{sec:lom}).
Repeat counts and timing boundaries differ across experiments and are
stated in each table. The large-scale eigenmode completion used a
different 128\,GB box and lacks mode-content verification; it cannot
establish equivalent-work throughput.

\paragraph{Limitations.}
(i) One geometry. The eigenmode results are on one mesh at first order;
higher-order and refined eigensolves are future work (a $p = 2$ Palace
configuration is committed for that purpose). The electrostatic branch
has a convergence study, and it shows the committed mesh far from
converged for $C_\Sigma$ (Section~\ref{sec:lom}). (ii) The lossless model
leaves readout ports open; resistive terminations (hence
complex-symmetric pencils and mode $Q$s) are out of scope. (iii) The
credential eigenpath remains the ungauged committed benchmark, and mode
classification still uses the Palace reference spectrum. The shipped
port term retains the $17.49$\,GHz port-sheet artifact; the rank-1
port-voltage formulation that removes it in a probe is planned and not
built. The committed artifacts still label the qubit mode as spurious. (iv) The CPU evidence depends on
request selection, field output, error estimation and implementation
version, and Palace solves about $17\%$ more unknowns from the same
input mesh; it supports no unconditional per-core or scaling advantage.
(v) GPU evidence covers
driven-solve correctness and an honest scaling negative in f32 on a
cube fixture family; no GPU eigensolve path exists. (vi) The
differentiable branch is electrostatic/LOM only (eigenmode/EPR is
roadmap) and scalar-$\varepsilon$ for the shape chain. Its real-device
result carries upper bounds and extrapolated estimates only, for the
polygonal geometry of the committed mesh and with the junction's
$5.5$\,fF outside $C_\Sigma$. The refined optimization runs in native
f64 with hand-derived element integrals, outside the Burn tape.

\section{Conclusion}
\label{sec:conclusion}

The discrete-adjoint layer supplies solver-derived material and
geometry gradients for electrostatic and driven FEM operators, and the
stationary capacitance derivative closes the geometry-to-charging-energy
loop. On the real device, a finite-difference-validated P2 shape
gradient on an adaptively refined mesh reaches the $89.9$\,fF anchor in
three accepted steps. Held-out refinement then shows what that design
is worth: its converged capacitance is rigorously below $89.02$\,fF and
is estimated at about $88.8$--$88.9$\,fF. A resolution-corrected design
converges below $90.08$\,fF, with an estimate band that contains the
anchor by construction. Only upper bounds are proven; no continuum
value is certified, and a two-sided bound is the next experiment. The
convergence study behind these statements is a result of its own: the
committed mesh overstates $C_\Sigma$ by at least $54.69$\,fF, the
converged device lies below the anchor and not above it, and the
coarse-mesh optimizations that lowered $C_\Sigma$ to reach the anchor
moved the wrong way while passing every self-consistency check.
Same-input-mesh agreement with Palace supplies a forward-solver
credential that now includes the transmon qubit mode, which earlier
versions wrongly reported as absent. The updated mode-verified timing
and driven AMS experiments delimit the computational cost without
claiming a general performance advantage. Eigenmode/EPR
derivatives remain future work.

\paragraph{Acknowledgments.}
geode-fem's development was AI-orchestrated (agent-driven implementation
with human direction and review); this benchmark program and paper
follow the same discipline, with every reported number traced to a
committed, human-auditable artifact.
% TODO(operator): finalize acknowledgment wording before submission.
% TODO(operator): whiteroom L1-L4 operator specification — decide
% public/citable (repo/DOI) vs acknowledged-not-cited (BRIEF coauthor note).

\bibliographystyle{plainnat}
\bibliography{refs}

\end{document}
